\documentclass[a4paper,fleqn]{cas-sc}

\usepackage[numbers]{natbib}
\usepackage[shortlabels]{enumitem}

\ExplSyntaxOn
\RenewDocumentCommand \emailauthor { m m }
{
  \int_gincr:N \g_ead_int
  \seq_gput_right:Nn \g_stm_ead_seq
    {
      { \ttfamily \tl_to_str:n { #1 } }
      \space(#2)%
    }
}
\RenewDocumentCommand \printorcid { } { }
\IfFileExists{thumbnails/cas-email.jpeg}{}{\bool_gset_true:N \g_stm_nologo_bool}

\cs_set:Npn \__reset_fig:
{
  \tl_set:Nx \l_fig_pos_tl { !htb }
  \tl_set:Nx \l_fig_cols_tl { 1 }
  \tl_set:Nn \l_fig_align_tl { \centering }
  \skip_set:Nn \l_fig_abovecap_skip { 3pt }
  \skip_set:Nn \l_fig_belowcap_skip { 3pt }
  \skip_set:Nn \l_fig_abovefig_skip { 3pt }
  \skip_set:Nn \l_fig_belowfig_skip { 3pt }
}
\cs_set:Npn \__reset_tbl:
{
  \tl_set:Nx \l_tbl_pos_tl { !htb }
  \tl_set:Nx \l_tbl_cols_tl { 1 }
  \tl_set:Nn \l_tbl_align_tl { \centering }
  \skip_set:Nn \l_tbl_abovecap_skip { 3pt }
  \skip_set:Nn \l_tbl_belowcap_skip { 2pt }
  \skip_set:Nn \l_tbl_abovetbl_skip { 3pt }
  \skip_set:Nn \l_tbl_belowtbl_skip { 3pt }
}

% Fix bug in cas-common.sty where \l_keys_key_tl is empty in modern expl3
\keys_define:nn { cas / fig }
{
  unknown .code:n =
    {
      \ifstrempty {#1}
        { \tl_set:NV \l_fig_pos_tl \l_keys_key_str }
        { \tl_set:Nn \l_fig_pos_tl { !htb } }
    }
}
\keys_define:nn { cas / tbl }
{
  unknown .code:n =
    {
      \ifstrempty {#1}
        { \tl_set:NV \l_tbl_pos_tl \l_keys_key_str }
        { \tl_set:Nn \l_tbl_pos_tl { !htb } }
    }
}
\ExplSyntaxOff

\raggedbottom
\makeatletter
\setlength{\@fptop}{0pt}
\setlength{\@fpbot}{0pt plus 1fil}
\makeatother

\renewcommand{\topfraction}{0.95}
\renewcommand{\bottomfraction}{0.95}
\renewcommand{\textfraction}{0.05}
\renewcommand{\floatpagefraction}{0.90}
\setcounter{topnumber}{10}
\setcounter{bottomnumber}{10}
\setcounter{totalnumber}{20}
\setcounter{dbltopnumber}{10}
\renewcommand{\dbltopfraction}{0.95}
\renewcommand{\dblfloatpagefraction}{0.90}
\setlength{\floatsep}{5pt plus 2pt minus 2pt}
\setlength{\intextsep}{5pt plus 2pt minus 2pt}
\setlength{\textfloatsep}{6pt plus 2pt minus 2pt}

\newcommand{\includesafeimage}[2][width=\textwidth]{%
    \IfFileExists{figures/#2}{\includegraphics[#1]{figures/#2}}{%
    \IfFileExists{paper/figures/#2}{\includegraphics[#1]{paper/figures/#2}}{%
    \IfFileExists{figs/#2}{\includegraphics[#1]{figs/#2}}{%
    \IfFileExists{#2}{\includegraphics[#1]{#2}}{%
    \IfFileExists{../paper/figures/#2}{\includegraphics[#1]{../paper/figures/#2}}{%
    \fbox{\parbox[c][3.2cm][c]{0.92\linewidth}{\centering \textbf{Figure Asset} \\[0.1cm] \texttt{\detokenize{#2}} \\[0.1cm] \footnotesize(Generate via \texttt{benchmarks/plot\_cns.py} and upload to Overleaf)}}%
    }}}}}%
}

\newdefinition{example}{Example}

\begin{document}
\let\WriteBookmarks\relax

\shorttitle{Exact SAT, CP, and ILP Approaches for SGP}
\shortauthors{D. H. Hieu et al.}

\title [mode = title]{Exact Approaches for the Social Golfer Problem: A Comparative Study of SAT Cardinality Encodings, Order Encoding, Constraint Programming, and Integer Linear Programming}

\author[1]{Do Huy Hieu}
\cormark[1]
\ead{24021470@vnu.edu.vn}

\author[1]{Pham Van Huong}
\cormark[1]
\ead{24021510@vnu.edu.vn}

\author[1]{Nguyen Minh Son}
\cormark[1]
\ead{24021614@vnu.edu.vn}

\affiliation[1]{organization={Faculty of Information Technology, VNU University of Engineering and Technology},
            addressline={144 Xuan Thuy, Cau Giay},
            city={Hanoi},
            country={Vietnam}}

\cortext[1]{Equal contribution. Group report for Course INT3011E -- Advanced Topics in Computer Science. Student IDs: 24021470 (D. H. Hieu), 24021510 (P. V. Huong), 24021614 (N. M. Son). Supervisor: Dr. To Van Khanh.}

\begin{abstract}
The Social Golfer Problem (SGP) asks whether $n = gs$ players can be scheduled into $g$ groups of size $s$ over $w$ weeks so that no two players end up in the same group more than once. In this report, we implement and compare four exact solving approaches in a shared codebase: a three-index one-hot SAT encoding tested with four \textit{Exactly}-$s$ encodings---New Sequential Counter (NSC), Sequential Counter, Totalizer, and Cardinality Network---using CaDiCaL~1.5.3 and Glucose~4.1; a slot-based Order Encoding; a Constraint Programming (CP) model in Google OR-Tools CP-SAT; and an Integer Linear Programming (ILP) model solved with Gurobi~14.0. Alongside an initial 2-second SAT screening across 528 configurations, we run a controlled single-thread, 60-second baseline without optional \texttt{SB1}/\texttt{SB2} (\texttt{SB=none}) over all 64 active reference instances (768 runs total), enforcing a 5{,}000{,}000-clause C3 lower-bound guard for CNF encodings and checking every returned schedule with an independent Universal Verifier. Across the 768 baseline runs, 330 reach verified SAT (covering 34 distinct instances), 84 are ruled UNSAT early by the theoretical partner bound, 20 hit the C3 clause guard, and 334 time out, with zero errors or status mismatches. Within 60\,s at \texttt{SB=none}, SAT \texttt{NSC}+Glucose solves the most instances (32, averaging 2.02\,s on the 30 instances also solved by \texttt{NSC}+CaDiCaL vs.\ 2.41\,s for CaDiCaL), OR-Tools CP-SAT solves 31 instances including \texttt{19-10-2} ($n=190$ in 36.10\,s), Gurobi ILP solves 24 instances, and Direct (NSC) achieves a $6.6\times$ lower mean runtime than Order Encoding across their 20 shared CaDiCaL cases (0.71\,s vs.\ 4.72\,s). We also run a 640-run symmetry-breaking ablation (\texttt{sb1} and \texttt{full}) across representative solvers and standalone CP-SAT runs on the two disabled instances with the largest estimated C3 clause counts (\texttt{11-11-12} and \texttt{13-13-14}). Within 60\,s, enabling \texttt{sb1} lets all six tested configurations prove 7 of the 8 remaining non-precheck UNSAT instances and speeds up Order Encoding by $3.1\times$--$4.8\times$. Under \texttt{full}, CP-SAT proves all 8 non-precheck UNSAT instances, speeds up by $2.8\times$, and decides 47/64 active instances on its own, while combining all tested configurations decides 50/64 active instances; on \texttt{11-11-12} and \texttt{13-13-14}, CP-SAT builds the models in 1.29\,s and 2.85\,s before timing out at 60\,s.
\end{abstract}

\begin{keywords}
Social Golfer Problem \sep SAT encoding \sep Cardinality constraints \sep New Sequential Counter \sep Order encoding \sep Constraint Programming \sep Integer Linear Programming
\end{keywords}

\maketitle

\section{Introduction}
\label{sec:introduction}

Timetabling and grouping problems appear often in discrete mathematics, operations research, and automated reasoning \cite{colbourn2006handbook,biere2021handbook}. Among them, the \textbf{Social Golfer Problem (SGP)}---listed as \textbf{CSPLib problem~010} \cite{harvey2001csplib}---is a well-known benchmark for testing exact constraint solvers. Given $g$ groups of $p$ players each (so $n = g \times p$ golfers in total) and $w$ weeks, the decision problem asks whether we can assign every player to one group each week so that no two players share a group more than once across all $w$ weeks. Whether this general decision problem (with no preassigned players) is in P or NP-complete is still open, whereas the completion version starting from a partial schedule is known to be NP-complete~\cite{triska2012tme,verplancke2026}.

The SGP is hard to solve for two main reasons: the search space grows exponentially, and any valid schedule has a huge number of symmetric copies \cite{harvey2005molr,triska2012tme}. With $A = \frac{(g \times p)!}{(p!)^g \, g!}$ ways to partition $n = g \times p$ players into $g$ unlabeled groups of size $p$ in one week, there are $A^w$ candidate schedules across $w$ ordered weeks before applying unique pairings (C3) or factoring out week and player permutations (dividing by $w!$ gives the unordered-week approximation $\frac{((g \times p)!)^w}{(p!)^{gw} \times (g!)^w \times w!}$ in~\cite{kieu2025nse} when weekly partitions are distinct). Even for moderate sizes, a solver can spend most of its time exploring equivalent symmetric assignments unless symmetry is reduced in the model or handled with extra symmetry-breaking constraints \cite{walsh2006symmetry}.

Besides being a classic benchmark, the SGP also captures practical scheduling tasks. In sports scheduling, it helps generate balanced tournament matchups and round-robin fixtures \cite{briskorn2008sports,henz2001basketball}. In experimental design, it balances how treatments are grouped into blocks to reduce bias \cite{colbourn2006handbook,stinson2004designs,montgomery2017design}. Similar grouping rules also show up when scheduling exam invigilators or rotating group assignments \cite{verplancke2026,triska2012tme} and assigning tables in competitive mahjong tournaments \cite{lester2021mahjong}.

Previous work on the SGP generally falls into four main approaches:
\begin{itemize}[noitemsep]
    \item \textbf{Constraint Programming (CP)}: Uses integer variables, global constraints such as \texttt{AllDifferent}, and static or dynamic symmetry breaking \cite{harvey2005molr,azevedo2006dynamic,liu2019sgp,ortools}.
    \item \textbf{Integer and Mixed-Integer Linear Programming (ILP/MIP)}: Uses binary variables for player-group-week assignments and linearizes pairwise meeting constraints for branch-and-cut solvers \cite{margot2002pruning,gurobi}.
    \item \textbf{Metaheuristics and Local Search}: Applies methods like tabu search or memetic algorithms to find schedules quickly on large inputs \cite{cotta2006memetic,dotu2007local}, though these incomplete methods cannot prove unsatisfiability (UNSAT).
    \item \textbf{Boolean Satisfiability (SAT)}: Converts the SGP into Conjunctive Normal Form (CNF) for CDCL solvers, moving from four-index Direct Encoding (DE) \cite{gent2005sat} and Triska--Musliu Encoding (TME) \cite{triska2012tme} to Set Constraint Encoding (SCE) \cite{lardeux2015sce} and the more compact three-index New SAT Encoding (NSE) with the New Sequential Counter (NSC) \cite{kieu2025nse}.
\end{itemize}

In standard SGP notation, instances are written as $g\text{-}p\text{-}w$ (groups, players per group, weeks). In our codebase and API, we write $(w, g, s)$ with group size $s = p$ and $n = g \times s$; for example, the small instance \texttt{2-2-3} corresponds to $\text{SGP}(w=3, g=2, s=2)$.

In this report, we design, implement, and compare SAT, Order Encoding, CP-SAT, and ILP solvers on the SGP under a shared test suite. Our main contributions are:
\begin{enumerate}[label=(\arabic*),noitemsep]
    \item \textbf{Comparing C2 Cardinality Encodings}: Keeping the three-index one-hot SAT variables with fixed group-assignment (C1) and unique-pairing (C3) clauses, we test four Group Size (C2) encodings---NSC \cite{kieu2025nse}, Sequential Counter \cite{sinz2005towards}, Totalizer \cite{bailleux2003totalizer}, and Cardinality Network \cite{asin2009cardnet}---on CaDiCaL~1.5.3 and Glucose~4.1, covering both formula sizes/2-second screening and a full 60-second benchmark.
    \item \textbf{Direct vs.\ Order Encoding and Symmetry-Breaking Ablation}: We compare three-index Direct SAT against slot-based Order Encoding~\cite{tamura2009order} under \texttt{SB=none}, and run a 640-run ablation contrasting \texttt{SB=none}, \texttt{sb1}, and \texttt{full} (\texttt{SB1}+\texttt{SB2})~\cite{kieu2025nse} across representative SAT, Order, CP-SAT, and ILP solvers.
    \item \textbf{Cross-Paradigm Comparison and Large-Scale Check}: We test the SAT encodings alongside Google OR-Tools CP-SAT (\texttt{AllDifferent}, \texttt{Inverse}, \texttt{AtMostOne}) and Gurobi~14.0 ILP under an identical single-thread, 60-second limit across 64 active instances, plus standalone CP-SAT runs on the two disabled instances with the largest estimated C3 clause counts (\texttt{11-11-12} and \texttt{13-13-14}).
    \item \textbf{Shared Reproducible Pipeline}: We build a common CLI (\texttt{main.py}), DIMACS CNF exporter, partner-bound precheck, standalone Universal Verifier (\texttt{src/\allowbreak verifier.py}), CSV~v1 benchmark runner (\texttt{benchmarks/\allowbreak run\_\allowbreak all.py}), and CNSPlots visualization script (\texttt{benchmarks/\allowbreak plot\_\allowbreak cns.py}).
\end{enumerate}

The rest of this report is organized as follows. Section~\ref{sec:formulation} defines the SGP, states the partner bound, and describes the CP and ILP models. Section~\ref{sec:sat_approach} covers the three-index SAT model, the four \textit{Exactly}-$s$ encodings, Order Encoding, symmetry breaking, and our verification pipeline. Section~\ref{sec:experiments} presents the experimental setup and benchmark results, and Section~\ref{sec:conclusion} summarizes our findings and future work.

\section{Problem Formulation and Exact Baseline Paradigms}
\label{sec:formulation}

\subsection{Formal Definition, Partner Bound, and Notation}
\label{sec:problem_def}

Let $P = \{0, 1, \ldots, n - 1\}$ denote the set of $n = g \times s$ players, where $g \ge 1$ is the number of groups per week, $s \ge 1$ is the group size, and $w \ge 1$ is the number of weeks. Using zero-based indices for weeks $k \in \{0, \ldots, w - 1\}$ and groups $j \in \{0, \ldots, g - 1\}$, a valid SGP schedule is a collection of subsets $\{G_{k,j} \subseteq P\}$ satisfying three conditions:
\begin{enumerate}[label=\textbf{(C\arabic*)},noitemsep]
    \item \textbf{Unique Group Assignment}: In every week $k$, the groups partition $P$: $\bigcup_{j=0}^{g-1} G_{k,j} = P$ and $G_{k,j_1} \cap G_{k,j_2} = \emptyset$ for $j_1 \neq j_2$.
    \item \textbf{Group Size}: Every group contains exactly $s$ players: $|G_{k,j}| = s$ for all $k \in \{0, \ldots, w - 1\}$ and $j \in \{0, \ldots, g - 1\}$.
    \item \textbf{Unique Pairings (Social Constraint)}: No two distinct players meet more than once across the entire schedule: $|G_{k_1,j_1} \cap G_{k_2,j_2}| \le 1$ whenever $(k_1, j_1) \neq (k_2, j_2)$.
\end{enumerate}

\noindent\textbf{Theoretical Partner Bound:} When $s > 1$, each player meets $s - 1$ different teammates per week out of the $n - 1 = gs - 1$ other players. Because pairs cannot repeat, any valid schedule must satisfy:
\begin{equation}
    w \le w_{\max} = \left\lfloor \frac{g s - 1}{s - 1} \right\rfloor \qquad (s > 1).
    \label{eq:partner_bound}
\end{equation}
When $s = 1$, each group has only one player, so there is no upper limit on $w$. Note that Equation~(\ref{eq:partner_bound}) is necessary but \emph{not sufficient}. For example, $\text{SGP}(w=2, g=2, s=3)$ satisfies $w = 2 \le \lfloor (6-1)/(3-1) \rfloor = 2$, yet it is still UNSAT because any 3-player group in week~2 must contain at least 2 players who were already together in week~1 (by the pigeonhole principle).

\begin{example}
Consider the instance \texttt{2-2-3} ($g = 2$, $p = s = 2$, $w = 3$, so $n = 4$ players). In 1-based display notation, a valid schedule achieving $w = w_{\max} = 3$ is: Week~1: $\{1, 2\}, \{3, 4\}$; Week~2: $\{1, 3\}, \{2, 4\}$; Week~3: $\{1, 4\}, \{2, 3\}$.
\end{example}

Table~\ref{tab:notation_map} summarizes how the notation of Kieu et al.~\cite{kieu2025nse} maps to our implementation API.

\begin{table}[pos=!htb]
\centering\footnotesize
\setlength{\tabcolsep}{4pt}
\renewcommand{\arraystretch}{0.85}
\caption{Notation mapping between the reference SGP paper \cite{kieu2025nse} and our implementation API.}
\label{tab:notation_map}
\begin{tabular}{lll}
\toprule
\textbf{Concept} & \textbf{Reference Paper \cite{kieu2025nse}} & \textbf{Project API / Code} \\
\midrule
Instance label / signature & $(g, p, w)$ (e.g., \texttt{2-2-3}) & $(w, g, s)$ with $s = p$ (e.g., \texttt{solve\_sgp(3, 2, 2)}) \\
Total players & $n = g \times p$ & $n = g \times s$ \\
Player, group, week indices & $i \in \{1..n\},\ k \in \{1..g\},\ l \in \{1..w\}$ & $i \in \{0..n{-}1\},\ j \in \{0..g{-}1\},\ k \in \{0..w{-}1\}$ \\
Group assignment variable & $G_{ikl} \in \{\text{TRUE}, \text{FALSE}\}$ & $\texttt{var}(i,j,k,g,w) = i\cdot g\cdot w + j\cdot w + k + 1$ \\
\bottomrule
\end{tabular}
\end{table}

\subsection{Constraint Programming (CP) Formulation}
\label{sec:cp_formulation}

In Constraint Programming \cite{harvey2005molr,ortools}, we model the SGP with integer variables and global constraints in \texttt{src/cp\_solver.py} using Google OR-Tools CP-SAT. Let $K = \{0, \ldots, w - 1\}$, $J = \{0, \ldots, g - 1\}$, $C = \{0, \ldots, s - 1\}$, and $P = \{0, \ldots, n - 1\}$.

\noindent\textbf{Variables:}
\begin{itemize}[noitemsep]
    \item $\text{slot}_{k,j,c} \in \{0, \ldots, n - 1\}$: the player assigned to slot $c \in C$ of group $j \in J$ in week $k \in K$.
    \item $\text{pos}_{k,i} \in \{0, \ldots, n - 1\}$: the flattened slot index $j \cdot s + c$ occupied by player $i \in P$ in week $k \in K$.
    \item $\text{group\_of}_{k,i} \in \{0, \ldots, g - 1\}$: the group index of player $i \in P$ in week $k \in K$.
    \item $m_{i_1,i_2,k} \in \{0, 1\}$: a Boolean indicator equal to $1$ if and only if players $i_1 < i_2$ are in the same group in week $k$.
\end{itemize}

\noindent\textbf{Constraints:} Each group always holds $s$ slot variables, so (C2) is satisfied automatically. Within each week $k \in K$, we add:
\begin{align}
    \text{slot}_{k,j,c} &< \text{slot}_{k,j,c+1}
      && j \in J,\ 0\le c<s-1, \label{eq:cp_order} \\
    &\texttt{AllDifferent}\big(\{\text{slot}_{k,j,c} : j \in J,\ c \in C\}\big)
      && \text{(C1)}, \label{eq:cp_alldiff} \\
    &\texttt{Inverse}\big(\text{slot}_{k,\cdot,\cdot},\ \text{pos}_{k,\cdot}\big)
      && \text{(Position channeling)}, \label{eq:cp_inverse} \\
    \text{group\_of}_{k,i} &= \lfloor \text{pos}_{k,i} / s \rfloor
      && i \in P. \label{eq:cp_div}
\end{align}
When $s = 1$, Equation~(\ref{eq:cp_div}) simply becomes $\text{group\_of}_{k,i} = \text{pos}_{k,i}$.

To enforce Unique Pairings (C3) when $w > 1$ and $s > 1$, instead of enumerating $O(w^2 g^2 n^2)$ combinations of player pairs, weeks, and groups, we track whether each pair meets in each week and add an \texttt{AtMostOne} constraint across weeks for every pair $0 \le i_1 < i_2 < n$:
\begin{align}
    m_{i_1,i_2,k} = 1 &\iff \big(\text{group\_of}_{k,i_1} = \text{group\_of}_{k,i_2}\big) && \forall k \in K, \label{eq:cp_reify} \\
    \sum_{k=0}^{w-1} m_{i_1,i_2,k} &\le 1 && \text{(\texttt{AddAtMostOne} across weeks)}. \label{eq:cp_social}
\end{align}
In total, the model uses $w n$ slot variables, $w n$ position variables, $w n$ group variables (for $s > 1$), and $w \binom{n}{2}$ meeting indicators when $w>1$ and $s>1$. If optional static symmetry breaking is enabled via \texttt{--sb}, we also fix the first week (\texttt{sb1}: $\text{slot}_{0,j,c} = j \cdot s + c$) and assign the first $\min(g, s)$ players to distinct groups in later weeks (\texttt{full}: $\text{group\_of}_{k,i} = i$ for $k \ge 1, 0 \le i < \min(g, s)$). Inside OR-Tools CP-SAT, \texttt{AllDifferent} narrows variable domains in each week, while built-in clause learning handles conflicts over the $m_{i_1,i_2,k}$ variables.

\subsection{Integer Linear Programming (ILP) Formulation}
\label{sec:ilp_formulation}

% [TV3] The ILP description follows the implemented model in src/ilp_solver.py.
In Integer Linear Programming \cite{gurobi}, the SGP is formulated using binary assignment variables $x_{i,j,k} \in \{0, 1\}$, where $x_{i,j,k} = 1$ if player $i \in P$ is assigned to group $j \in J$ in week $k \in K$, and $0$ otherwise. To linearize the quadratic co-occurrence condition $x_{i_1,j,k} \cdot x_{i_2,j,k}$, auxiliary continuous variables $y_{i_1,i_2,j,k} \in [0, 1]$ are introduced for each pair $0 \le i_1 < i_2 < n$, group $j \in J$, and week $k \in K$:
\begin{align}
    \text{find} \quad & x_{i,j,k} \in \{0, 1\},\ \ 0 \le y_{i_1,i_2,j,k} \le 1 \label{eq:ilp_vars} \\
    \text{s.t.} \quad & \sum_{j=0}^{g-1} x_{i,j,k} = 1 && \forall i \in P,\ k \in K \quad \text{(C1)}, \label{eq:ilp_c1} \\
                      & \sum_{i=0}^{n-1} x_{i,j,k} = s && \forall j \in J,\ k \in K \quad \text{(C2)}, \label{eq:ilp_c2} \\
                      & y_{i_1,i_2,j,k} \ge x_{i_1,j,k} + x_{i_2,j,k} - 1
                          && i_1<i_2,\ j \in J,\ k \in K, \label{eq:ilp_lin}\\
                      & y_{i_1,i_2,j,k} \le x_{i_1,j,k},\quad y_{i_1,i_2,j,k} \le x_{i_2,j,k}
                          && i_1<i_2,\ j \in J,\ k \in K, \label{eq:ilp_lin_ub}\\
                      & \sum_{k=0}^{w-1} \sum_{j=0}^{g-1} y_{i_1,i_2,j,k} \le 1 && \forall i_1 < i_2 \quad \text{(C3)}. \label{eq:ilp_c3}
\end{align}
The model has no objective: Gurobi is asked only for a feasible schedule. The continuous auxiliary variables are exact for this formulation. Together with $0 \le y_{i_1,i_2,j,k} \le 1$, the three McCormick linearization inequalities~\eqref{eq:ilp_lin}--\eqref{eq:ilp_lin_ub} uniquely force $y_{i_1,i_2,j,k} = x_{i_1,j,k} x_{i_2,j,k} \in \{0, 1\}$ for every binary assignment $x \in \{0, 1\}^{ngw}$: inequality~\eqref{eq:ilp_lin} forces $y_{i_1,i_2,j,k}=1$ when both players share group $j$ in week $k$, while the upper bounds~\eqref{eq:ilp_lin_ub} force $y_{i_1,i_2,j,k}=0$ otherwise (and tighten the LP relaxation). Hence explicit integrality constraints on $y$ are unnecessary and the branch-and-cut search never branches on these variables.

\paragraph{Model size.} Without symmetry-breaking constraints, the model contains $ngw$ binary and $\binom{n}{2}gw$ continuous variables, and
\begin{equation}
    C_{\mathrm{ILP}} = nw + gw + 3\binom{n}{2}gw + \binom{n}{2}
    \label{eq:ilp_size}
\end{equation}
linear constraints: the $nw$ assignment equalities~\eqref{eq:ilp_c1}, the $gw$ group-size equalities~\eqref{eq:ilp_c2}, the three McCormick inequalities~\eqref{eq:ilp_lin}--\eqref{eq:ilp_lin_ub} for every $y$, and the $\binom{n}{2}$ meeting constraints~\eqref{eq:ilp_c3}. For instance, \texttt{3-2-2} ($w=2$, $g=3$, $s=2$, $n=6$) gives $36$ binary variables, $90$ continuous variables, and $303$ constraints, which agrees with the model reported by Gurobi.

\paragraph{Symmetry breaking.} The optional \texttt{--sb} modes add the unit constraints of Section~\ref{sec:symmetry_breaking} to the same model: \texttt{sb1} fixes the first week, and \texttt{full} also fixes the first $\min(g,s)$ players in later weeks. These rows are linear equalities $x_{i,j,k}=1$ and are counted separately from~\eqref{eq:ilp_size}.

\paragraph{Precheck, statuses, and timing.} The partner bound~\eqref{eq:partner_bound} is evaluated before a model is built, so a bound-violating instance returns \texttt{UNSAT} without a Gurobi call. A single Gurobi environment is created with output disabled and reused; dual reductions are disabled so that an infeasible model receives a definitive \texttt{INFEASIBLE} status instead of \texttt{INF\_OR\_UNBD}. The time limit and the thread count are set on the model. An optimal solution is decoded and passed through the shared verifier. \texttt{INFEASIBLE} maps to \texttt{UNSAT}; when the time limit is reached, a verified incumbent is retained as \texttt{SAT} and otherwise the run is reported as \texttt{TIMEOUT}; a missing license or an exceeded size limit is reported as \texttt{UNAVAILABLE}. Model-construction time, Gurobi's reported solve time, process CPU time, and whole-run wall time are recorded separately.

\section{SAT-Based Approaches}
\label{sec:sat_approach}

\subsection{Three-Index One-Hot Representation}
\label{sec:sat_model}

We use the player--group--week representation of the New SAT Encoding (NSE)~\cite{kieu2025nse}. The original four-index Direct Encoding (DE) also models positions within groups~\cite{gent2005sat}. In this study, \emph{direct} in the comparison with Order Encoding refers to the three-index assignment model.
We use Boolean variables
\begin{equation}
 x_{i,j,t}=1
 \quad\Longleftrightarrow\quad
 \text{player $i$ belongs to group $j$ in week $t$},
 \label{eq:sgp-direct-variable}
\end{equation}
where $1\leq i\leq n$, $1\leq j\leq g$, and $1\leq t\leq w$.
Thus the model contains $V_{\mathrm{primary}}=ngw$ primary variables.
This subsection uses one-based indices; the Python API uses
zero-based player, group, and week indices. For zero-based Python indices $(i_0,j_0,t_0)$,
the DIMACS identifier is $\texttt{var}(i_0,j_0,t_0,g,w)=i_0 gw+j_0 w+t_0+1$. This bijection reserves identifiers
$1,\ldots,ngw$ for primary variables. Auxiliary variables are allocated
above this range, with a monotonically increasing identifier shared by
successive cardinality constraints.

\paragraph{C1: one group per player and week.}
Each player-week pair must select exactly one group. The implementation
combines one at-least-one clause with the pairwise at-most-one clauses:
\begin{align}
 &\bigvee_{j=1}^{g} x_{i,j,t},
 &&1\leq i\leq n,\quad 1\leq t\leq w,
 \label{eq:sgp-c1-alo}\\
 &\neg x_{i,j,t}\lor\neg x_{i,j',t},
 &&1\leq i\leq n,\quad 1\leq t\leq w,\quad 1\leq j<j'\leq g.
 \label{eq:sgp-c1-amo}
\end{align}
The pairwise choice is fixed in all four C2-ablation configurations. It
introduces no auxiliary variables and emits exactly
\begin{equation}
 C_1=nw\left(1+\binom{g}{2}\right)
 \label{eq:sgp-c1-size}
\end{equation}
clauses, including the unit at-least-one clauses when $g=1$.

\paragraph{C2: fixed group size.}
Every group contains exactly $s$ players:
\begin{equation}
 \sum_{i=1}^{n}x_{i,j,t}=s,
 \qquad 1\leq j\leq g,\quad 1\leq t\leq w.
 \label{eq:sgp-c2}
\end{equation}
Each C2 constraint is an equality with target $s$. Section~\ref{sec:c2_encodings} describes the four
alternative CNF representations of~\eqref{eq:sgp-c2}.

\paragraph{C3: no repeated pair.}
For each player pair $i<i'$, distinct weeks $t<t'$, and groups $j,j'$, we
add the clause
\begin{equation}
 \neg x_{i,j,t}\lor\neg x_{i',j,t}\lor
 \neg x_{i,j',t'}\lor\neg x_{i',j',t'}.
 \label{eq:sgp-c3}
\end{equation}
The groups $j$ and $j'$ may coincide. A clause is violated precisely
when the two players share a group in both selected weeks. Consequently,
these clauses forbid every repeated meeting and allow pairs that never
meet. They introduce no auxiliary variables and their exact count is
\begin{equation}
 C_3=\binom{w}{2}g^2\binom{n}{2}.
 \label{eq:sgp-c3-size}
\end{equation}
The pairwise form is simple but grows quadratically in both the number
of players and the number of weeks. Its cost must therefore be included
when interpreting total CNF size.

\paragraph{Correctness and decoding.}
For an assignment satisfying C1, each player has one group in each week.
C2 gives all groups the required size, and C3 prevents repeated player
pairs. Reading the true primary variables therefore produces a valid
SGP schedule. Conversely, a valid schedule satisfies C1--C3, and each
correct cardinality encoding admits an auxiliary extension of its C2
assignments. Hence the conjunction is satisfiable exactly when a valid
schedule exists. Decoding ignores auxiliary variables. The runner checks that the Boolean model satisfies the CNF, then verifies
weekly coverage, group sizes and pair uniqueness in the decoded schedule.
This semantic argument is distinct from the finite regression tests:
exhaustive checks on small domains can reveal implementation errors but
do not by themselves prove correctness for every parameter value.

\subsection{Exactly-$s$ Cardinality Encodings for Group Size}
\label{sec:c2_encodings}

For one fixed group and week, let $a_i=x_{i,j,t}$ and let $k=s$.
We compare four implementations of $\sum_{i=1}^{n}a_i=k$: New
Sequential Counter (\texttt{nsc}), Sequential Counter
(\texttt{seqcounter}), Totalizer (\texttt{totalizer}), and Cardinality
Network (\texttt{cardnetwrk}). These are four configurations of the same
compact three-index model. We hold C1 and C3 fixed and vary only the CNF for C2.

\paragraph{New Sequential Counter.}
The NSC implementation follows the counter clauses of
Kieu et al.~\cite{kieu2025nse}. For $1\leq k<n$, it introduces
\begin{equation}
 R_{i,h}\ \Longleftrightarrow\
 \sum_{u=1}^{i}a_u\geq h,
 \qquad
 1\leq i\leq n-1,\quad 1\leq h\leq\min(i,k).
 \label{eq:nsc-meaning}
\end{equation}
The group and week indices on $R$ are omitted here; each group-week pair
has its own auxiliary variables. The initialization clauses are
\begin{align}
 &\neg a_i\lor R_{i,1}, &&1\leq i\leq n-1,
 \label{eq:nsc-init-one}\\
 &a_i\lor\neg R_{i,i}, &&1\leq i\leq k.
 \label{eq:nsc-init-diagonal}
\end{align}
For $2\leq i\leq n-1$, the update clauses are
\begin{align}
 &\neg R_{i-1,h}\lor R_{i,h},
 &&1\leq h\leq\min(i-1,k),
 \label{eq:nsc-forward}\\
 &\neg a_i\lor\neg R_{i-1,h-1}\lor R_{i,h},
 &&2\leq h\leq\min(i,k),
 \label{eq:nsc-increment}\\
 &a_i\lor R_{i-1,h}\lor\neg R_{i,h},
 &&1\leq h\leq\min(i-1,k),
 \label{eq:nsc-false-input}\\
 &R_{i-1,h-1}\lor\neg R_{i,h},
 &&2\leq h\leq\min(i,k).
 \label{eq:nsc-reverse}
\end{align}
Together these clauses implement the threshold recurrence
$R_{i,h}=R_{i-1,h}\lor(a_i\land R_{i-1,h-1})$, with the unavailable
boundary thresholds interpreted as constants. The initialization handles
those boundaries without allocating constant variables. Finally, the
at-least-$k$ condition is represented by
\begin{equation}
 (R_{n-1,k}\lor a_n)\land
 (R_{n-1,k}\lor R_{n-1,k-1}),
 \label{eq:nsc-lower-bound}
\end{equation}
where the second clause is omitted for $k=1$. The at-most-$k$ condition
uses
\begin{equation}
 \neg a_i\lor\neg R_{i-1,k},\qquad k+1\leq i\leq n.
 \label{eq:nsc-upper-bound}
\end{equation}
By the threshold meaning of $R$, \eqref{eq:nsc-lower-bound} states that
at least $k$ inputs are true. If more than $k$ inputs were true, the
$(k+1)$-st true input would violate~\eqref{eq:nsc-upper-bound}. Conversely,
any assignment with exactly $k$ true inputs extends to the actual prefix
thresholds and satisfies all these clauses. This establishes the required
equisatisfiability with Exactly-$k$. The cases $k=0$ and $k=n$ are emitted
as $n$ negative or positive unit clauses, respectively, without counters.

The number of counter variables for a nondegenerate constraint is
\begin{equation}
 A(n,k)=\sum_{i=1}^{n-1}\min(i,k)
       =\frac{k(2n-k-1)}{2}.
 \label{eq:nsc-aux-count}
\end{equation}
Counting every clause appended by the implementation, including clauses
that could be removed by subsequent simplification, gives
\begin{equation}
 C_{\mathrm{NSC}}(n,k)
 =4A(n,k)-2k+2+\begin{cases}0,&k=1,\\1,&k\geq2.\end{cases}
 \label{eq:nsc-clause-count}
\end{equation}
For example, the four update families contribute
$2[A(n,k)-k]+2[A(n,k)-n+1]$ clauses. Adding the initialization and final
conditions yields~\eqref{eq:nsc-clause-count}. These are counts of this
implementation, rather than a claim that all formulations called NSC have
the same size.

\paragraph{Distinguishing source summary and emitted counts.}
Table~1 of Kieu et al.~\cite{kieu2025nse} lists the NSC clause expression
$2nk+n-3k-k^2$, whereas explicit generation of its Eq.~(2)--(9), with
the stated index ranges and CNF expansion of Eq.~(8), gives
\eqref{eq:nsc-clause-count} in this implementation. These expressions
are not identical: for $n=4$, $k=2$, both approaches to the auxiliary
count give five variables, but the summary expression gives ten
clauses and the explicit clause families emit nineteen. We report the
emitted counts for all four encodings under the same counting convention.

\paragraph{Sequential Counter.}
The sequential-counter family encodes a bound through prefix counters
following Sinz~\cite{sinz2005towards}. We use PySAT's \texttt{seqcounter}
implementation. The library source documents an irredundant variant,
so the original rectangular-counter formula $k(n-1)$ is not substituted
for the number actually generated~\cite{pysat-seqcounter-source}.

\paragraph{Totalizer.}
The Totalizer recursively combines unary representations of counts in a
binary tree~\cite{bailleux2003totalizer}. A merge relates thresholds in the two
children to thresholds in their parent. The project calls
\texttt{CardEnc.equals} with \texttt{EncType.totalizer}; it does not
share a counter tree between the lower and upper bounds.

\paragraph{Cardinality Network.}
Cardinality networks use comparator-based selection and merge structures
to expose the output threshold needed by a cardinality
bound~\cite{asin2009cardnet}. We select PySAT's \texttt{cardnetwrk} implementation.
Its generated size includes padding variables and clauses required by
the concrete network construction. Smaller CNF size alone does not imply
shorter SAT solving time.

\paragraph{Both sides of equality.}
For a library encoder $e$, \texttt{CardEnc.equals}
constructs both bounds~\cite{pysat-card-docs}:
\begin{equation}
 \sum_i a_i\leq k
 \quad\land\quad
 \sum_i \neg a_i\leq n-k.
 \label{eq:equals-two-bounds}
\end{equation}
Let $a_e(n,q)$ and $c_e(n,q)$ denote the actual auxiliary-variable and
clause counts of the selected library encoder for one AtMost-$q$
constraint. With disjoint auxiliary identifiers, one Exactly-$k$
constraint therefore contributes
\begin{align}
 A_e(n,k)&=a_e(n,k)+a_e(n,n-k),
 \label{eq:equals-aux-count}\\
 C_e(n,k)&=c_e(n,k)+c_e(n,n-k).
 \label{eq:equals-clause-count}
\end{align}
The common shortcuts are $a_e(n,0)=a_e(n,n-1)=a_e(n,n)=0$ and
$c_e(n,0)=n$, $c_e(n,n-1)=1$, $c_e(n,n)=0$, for $n\geq2$.
For $n=1$, the $q=0$ and $q=n$ cases suffice. The project's explicit
Exactly-$0$ and Exactly-$n$ shortcuts both use $n$ unit clauses.

\begin{table}[pos=!htb]
 \centering
 \caption{Cost per group-week C2 constraint. All rows encode equality;
 bounds are for emitted CNF before solver preprocessing. NSC formulas
 assume $1\leq k<n$. For the library rows, exact counts use
 \eqref{eq:equals-aux-count}--\eqref{eq:equals-clause-count}.}
 \label{tab:cardinality-complexity}
 \small
 \begin{tabular}{@{}lll@{}}
  \toprule
  Encoding & Auxiliary variables & Clauses \\
  \midrule
  NSC & $A(n,k)=k(2n-k-1)/2$ & $4A(n,k)-2k+2+\mathbf{1}_{k\geq2}$ \\
  Sequential Counter & $A_{\mathrm{SC}}(n,k)$ & $C_{\mathrm{SC}}(n,k)$ \\
  Totalizer & $O(n\log n)$ & $O(n^2)$ \\
  Cardinality Network & $O(n\log^2 n)$ & $O(n\log^2 n)$ \\
  \bottomrule
 \end{tabular}
\end{table}

For PySAT's irredundant sequential-counter implementation, the exact
nontrivial single-bound counts, for $1\leq q\leq n-2$, are
\begin{equation}
 a_{\mathrm{SC}}(n,q)=q(n-q),\qquad
 c_{\mathrm{SC}}(n,q)=(n-q)(2q+1)-q,
 \label{eq:sc-implementation-count}
\end{equation}
which determine $A_{\mathrm{SC}}(n,k)$ and $C_{\mathrm{SC}}(n,k)$ in
Table~\ref{tab:cardinality-complexity} via~\eqref{eq:equals-aux-count}--\eqref{eq:equals-clause-count}.
They follow by counting the allocated counter cells and generated clauses
in its nested loops~\cite{pysat-seqcounter-source}. The full unary
Totalizer and the comparator-network recurrences justify the conservative
worst-case bounds in Table~\ref{tab:cardinality-complexity}; the bounds
allow both sides of~\eqref{eq:equals-two-bounds} and implementation
padding~\cite{pysat-totalizer-source,pysat-network-source}. Bound-specific
shortcuts can make small instances substantially smaller. Accordingly,
we record the PySAT version and exact emitted variable and clause counts
directly from \texttt{src/sat\_direct\_encoders.py} rather than using
asymptotic estimates as measurements.

For every fully constructed three-index SGP formula without extra
symmetry clauses,
the total counts are
\begin{align}
 V_{\mathrm{SGP},e}&=ngw+gwA_e(n,s),
 \label{eq:sgp-total-vars}\\
 C_{\mathrm{SGP},e}&=nw\left(1+\binom{g}{2}\right)
      +gwC_e(n,s)+\binom{w}{2}g^2\binom{n}{2}.
 \label{eq:sgp-total-clauses}
\end{align}
In these equations $A_e(n,s)$ and $C_e(n,s)$ are the per-C2 counts; for
NSC with $s<n$ they are supplied by~\eqref{eq:nsc-aux-count} and
\eqref{eq:nsc-clause-count}. When $s=n$, $A_e(n,s)=0$ and $C_e(n,s)=n$.
The optional \texttt{symmetry\_clauses} argument to this encoder accepts
only primary literals. The added clauses are counted separately.

\subsection{Order Encoding}
\label{sec:order_encoding}

Order Encoding \cite{tamura2009order} replaces the one-hot membership
variables of Section~\ref{sec:sat_model} by an integer slot value and a
unary representation of that value. The model is implemented in
\path{src/order_encoder.py}; it is a slot-based model and shares the
$(w,g,s)$ API, the partner-bound precheck, and the standalone verifier
with the direct model.

\paragraph{Variables.}
Let $K=\{0,\ldots,w-1\}$, $J=\{0,\ldots,g-1\}$,
$C=\{0,\ldots,s-1\}$, and $V=\{1,\ldots,n\}$ with $n=gs$. Each group has
$s$ ordered slots. The integer variable $P_{k,j,c}\in V$ is the player
in slot $c$ of group $j$ in week $k$. For $n>1$, it is represented by
monotone Boolean variables
\begin{equation}
 B_{k,j,c,v}\ \Longleftrightarrow\ (P_{k,j,c}\le v),
 \qquad 1\le v\le n-1.
 \label{eq:order-meaning}
\end{equation}
The value $n$ is represented by all $B_{k,j,c,v}$ false and the value
$1$ by all of them true. The model contains $wgs(n-1)$ order variables;
the case $n=1$ is trivial and emits no clauses.

\paragraph{Order axiom.}
Every threshold implies the next one:
\begin{equation}
 \neg B_{k,j,c,v}\lor B_{k,j,c,v+1},
 \qquad 1\le v\le n-2.
 \label{eq:order-chain}
\end{equation}
This is the only constraint that links the $n-1$ thresholds of one slot.

\paragraph{Strict slot order.}
Within a group, slots are kept in increasing player order. For
$X=P_{k,j,c}$ and $Y=P_{k,j,c+1}$ with $c<s-1$, the clauses
\begin{align}
 &\neg B_{k,j,c+1,1},\qquad B_{k,j,c,n-1}, \label{eq:order-strict-bounds}\\
 &\neg B_{k,j,c+1,v}\lor B_{k,j,c,v-1},
 \qquad 2\le v\le n-1 \label{eq:order-strict}
\end{align}
encode $X<Y$: the first two clauses give $Y\ge 2$ and $X\le n-1$, and
\eqref{eq:order-strict} encodes
$(Y\le v)\Rightarrow(X\le v-1)$. Sorting the players of a group is
always possible, so this canonical order does not remove solutions.

\paragraph{Weekly assignment.}
C2 holds by construction, because each group has exactly $s$ slots. For
C1, every player must occupy exactly one slot in each week. The encoding
counts, for each threshold $v$, how many slots satisfy
$P_{k,j,c}\le v$:
\begin{equation}
 \sum_{j\in J}\sum_{c\in C} B_{k,j,c,v} = v,
 \qquad 1\le v\le n-1.
 \label{eq:order-alldiff}
\end{equation}
Each equality is compiled by PySAT's sequential counter
(\texttt{CardEnc.equals} with \texttt{EncType.seqcounter}); these are the
only auxiliary variables of the model. Let $c_v$ be the number of slots
with value at most $v$. Equation~\eqref{eq:order-alldiff} states
$c_v=v$, so $c_v-c_{v-1}=1$ for every $v$, with $c_0=0$ and $c_n=n$.
Hence the $gs$ slots carry the values $1,\ldots,n$ exactly once, which
is C1 for the week.

\paragraph{Membership channeling.}
C3 is expressed over membership indicators $e_{k,i,j}$, where
player $i$ appearing in group $j$ in week $k$ forces $e_{k,i,j}=1$. The
slot value implies membership:
\begin{equation}
 P_{k,j,c}=i\ \longrightarrow\ e_{k,i,j}.
 \label{eq:order-channel}
\end{equation}
At the clause level this is
\begin{align}
 &\neg B_{k,j,c,1}\lor e_{k,1,j}, \qquad
   B_{k,j,c,n-1}\lor e_{k,n,j}, \label{eq:order-channel-ends}\\
 &\neg B_{k,j,c,v}\lor B_{k,j,c,v-1}\lor e_{k,v,j},
   \qquad 2\le v\le n-1, \label{eq:order-channel-mid}
\end{align}
for the cases $i=1$, $i=n$, and $2\le i\le n-1$, respectively.
Because slots increase strictly, slot $c$ can only hold values
$c+1,\ldots,n-s+c+1$, so only those clauses are emitted. The implication
is one-directional, which is sufficient: C3 uses $e$ only negatively,
and every actual membership forces $e=1$, so unused indicators can
always be set to $0$. The $e$ variables contribute $wng$ bits.

\paragraph{No repeated pair.}
For $i_1<i_2$, weeks $k_1<k_2$, and groups $j_1,j_2$ (which may
coincide),
\begin{equation}
 \neg e_{k_1,i_1,j_1}\lor\neg e_{k_1,i_2,j_1}\lor
 \neg e_{k_2,i_1,j_2}\lor\neg e_{k_2,i_2,j_2}.
 \label{eq:order-c3}
\end{equation}
The clause forbids exactly the case in which the two players share group
$j_1$ in week $k_1$ and group $j_2$ in week $k_2$.

\paragraph{Size.}
For $n>1$ and without symmetry clauses, the model uses
\begin{equation}
 V_{\mathrm{ord}} = wgs(n-1)+wng
 \label{eq:order-vars}
\end{equation}
variables before the cardinality auxiliaries, and the non-cardinality clause groups
contribute
\begin{equation}
 C_{\mathrm{ord}} = wgs(n-2)+wg(s-1)n+wgs(n-s+1)
   +\binom{w}{2}g^2\binom{n}{2}
 \label{eq:order-clauses}
\end{equation}
clauses from~\eqref{eq:order-chain}, \eqref{eq:order-strict-bounds}--\eqref{eq:order-strict}, the
channeling clauses~\eqref{eq:order-channel-ends}--\eqref{eq:order-channel-mid}, and~\eqref{eq:order-c3}, respectively. In addition to $V_{\mathrm{ord}}$ and $C_{\mathrm{ord}}$, the
$w(n-1)$ threshold equalities~\eqref{eq:order-alldiff} contribute $w\sum_{v=1}^{n-1} A_{\mathrm{SC}}(n,v)$ sequential-counter auxiliary variables and $w\sum_{v=1}^{n-1} C_{\mathrm{SC}}(n,v)$ cardinality clauses (following~\eqref{eq:equals-aux-count}--\eqref{eq:equals-clause-count} and~\eqref{eq:sc-implementation-count}) to the total emitted CNF; on large instances these sequential-counter clauses can even exceed $C_{\mathrm{ord}}$, so the implementation reports the measured totals directly rather than estimating them. The SB1 variant of
Section~\ref{sec:symmetry_breaking} adds $gs(n-1)$ further unit clauses.

\paragraph{Correctness and decoding.}
Given \eqref{eq:order-chain}, the truth values of
$B_{k,j,c,\cdot}$ define a unique value $P_{k,j,c}$: the smallest $v$
with $B_{k,j,c,v}$ true, or $n$ when no threshold is true.
Equation~\eqref{eq:order-alldiff} makes the values a permutation of $V$
in each week and the strict-order clauses fix the slot order, so C1 and
C2 hold. Equation~\eqref{eq:order-c3}, together with
\eqref{eq:order-channel}, forbids repeated meetings. Conversely, any
valid schedule becomes a model of the encoding after sorting each
group, setting the thresholds, and setting $e$ on the memberships. The
decoder therefore returns the schedule from the slot values, and the
shared verifier checks C1--C3 again.

\subsection{Symmetry Breaking Constraints}
\label{sec:symmetry_breaking}

Every SGP instance is invariant under three families of label
permutations: a global relabeling of players, an independent relabeling
of groups within each week, and a permutation of the weeks. Isomorphic
copies of the same schedule can dominate the search space, and a CDCL
solver may explore many of them. Static symmetry breaking (SB) removes
part of this redundancy with unit clauses over primary
variables~\cite{kieu2025nse,walsh2006symmetry}. The implementation in
\path{src/symmetry.py} generates these clauses for the SAT encoders; the
CP-SAT and ILP models add the same restrictions directly.

\paragraph{SB1: canonical first week.}
For $j\in\{1,\ldots,g\}$ and $c\in\{1,\ldots,s\}$,
\begin{equation}
 x_{(j-1)s+c,\,j,\,1}=1.
 \label{eq:sb1}
\end{equation}
Week~1 is partitioned into the consecutive blocks
$\{1,\ldots,s\}$, $\{s+1,\ldots,2s\}$, \ldots, and block $j$ is assigned
to group $j$. Equation~\eqref{eq:sb1} uses $gs$ unit clauses and fixes
the player and group labels of the first week.

\paragraph{SB2: fixed players in later weeks.}
For every week $2\le t\le w$ and every player $1\le i\le\min(g,s)$,
\begin{equation}
 x_{i,i,t}=1.
 \label{eq:sb2}
\end{equation}
SB2 pins player $i$ to group $i$ in all weeks after the first and adds
$(w-1)\min(g,s)$ unit clauses. The implementation offers three modes:
\texttt{none} adds no SB clause, \texttt{sb1} adds \eqref{eq:sb1}, and
\texttt{full} adds \eqref{eq:sb1} together with \eqref{eq:sb2}.

\paragraph{Preservation under SB1.}
Let $S$ be a valid schedule with first-week groups $G_1,\ldots,G_g$.
Define a player permutation $\pi$ that maps the $s$ players of $G_j$ to
$\{(j-1)s+1,\ldots,js\}$, and apply $\pi$ to every week of $S$ while
keeping the group labels. The image is a valid schedule, because a
global player permutation preserves group sizes and meetings, and its
first week is exactly the block partition required by \eqref{eq:sb1}.
Hence SB1 removes only isomorphic copies, and appending \eqref{eq:sb1}
cannot turn a satisfiable instance into UNSAT.

\paragraph{Preservation under SB1+SB2.}
After \eqref{eq:sb1}, players $1,\ldots,\min(g,s)$ all belong to group
1 in week~1, since $\min(g,s)\le s$. They meet in that week, so C3
places them in pairwise distinct groups in every week $t\ge2$. For each
such week, let $\sigma_t$ be a permutation of the group labels that maps
the group of player $i$ to group $i$ for $i\le\min(g,s)$, and relabel
the groups of week $t$ by $\sigma_t$. This leaves the first week, C1,
C2, and the set of meetings unchanged, and produces $x_{i,i,t}=1$ for
every required pair. Therefore every feasible schedule has a relabeled
copy satisfying \eqref{eq:sb1} and \eqref{eq:sb2}, so \texttt{full}
preserves satisfiability.

\paragraph{Model-specific support.}
For the direct encoder, SB clauses are appended as unit clauses and
counted separately from the C2 encodings. CP-SAT adds
\eqref{eq:sb1}--\eqref{eq:sb2} as domain constraints, and the ILP model
adds the same rows as linear equalities. In the current implementation,
Order Encoding supports only \texttt{sb1} and rejects \texttt{--sb full}:
although pinning an arbitrary player to a group in a slot model generally
requires a disjunction over all $s$ slots, for the specific \texttt{SB2}
prefix $i\in\{1,\ldots,\min(g,s)\}$ in week $k\ge 1$ (using the
zero-based week and group indices of Section~\ref{sec:order_encoding}),
\texttt{SB2} pins all smaller players $1,\ldots,i-1$ to groups
$0,\ldots,i-2$, so C1 excludes them from group $i-1$ and strict slot
ordering~\eqref{eq:order-strict-bounds}--\eqref{eq:order-strict} forces
player $i$ into slot $c=0$ of group $i-1$ (i.e., $P_{k,i-1,0}=i$);
consequently, \texttt{SB2} is representable by unit clauses over
$B_{k,i-1,0,v}$ (at most two per pinned player and week via~\eqref{eq:order-chain},
or $n-1$ matching the full threshold pattern of \texttt{sb1}), and adding it
to \texttt{src/order\_encoder.py} is left as straightforward future work.
For \texttt{sb1}, the order model fixes the threshold pattern of week~0
to the value $js+c+1$ in slot $c$ of group $j$, using the zero-based
$j,c$ of Section~\ref{sec:order_encoding}; this is the slot-wise form of~\eqref{eq:sb1}.

\paragraph{Experimental use.}
The benchmark driver keeps \texttt{sb=none} as the default and records
the selected mode in the run metadata, so symmetry-breaking variants can
be compared on the same instances and encoders.

\subsection{Unified Solving Pipeline, DIMACS Export, and Universal Verification}
\label{sec:pipeline}

Our CLI (\texttt{main.py}) and benchmark runner (\texttt{benchmarks/run\_all.py}) run all four methods---\texttt{sat}, \texttt{order}, \texttt{cp}, and \texttt{ilp}---through the same pipeline and output format:
\begin{enumerate}[label=(\arabic*),noitemsep]
    \item \textbf{Partner-Bound Precheck}: All four solvers first check Equation~\eqref{eq:partner_bound} when $s > 1$. Any instance with $w > w_{\max}$ stops right away with \texttt{UNSAT\_PRECHECK} before building a model.
    \item \textbf{Model Construction \& DIMACS Export}: The SAT and Order solvers generate CNF formulas (with DIMACS file export supported via \texttt{--export-cnf} for \texttt{--method sat}), while CP-SAT and Gurobi build integer and linear models. Building time (\texttt{encode\_time\_seconds}), variable count (\texttt{num\_vars}), and clause or constraint count (\texttt{num\_clauses} / \texttt{num\_constraints}) are logged separately.
    \item \textbf{Isolated Runs \& Timing}: Each test case runs in its own worker process, initializing solver libraries and Gurobi before starting the total timeout clock (\texttt{total\_time\_seconds}), while also saving the solver-only time (\texttt{solver\_wall\_time\_seconds}). Results are written to a locked CSV~v1 file alongside a \texttt{.meta.json} settings file.
    \item \textbf{Universal Solution Verification}: Every \texttt{SAT} solution from any of the four solvers is converted to the 0-based format \texttt{schedule[week]}\allowbreak\texttt{[group]} and checked by \texttt{verify\_schedule} in \texttt{src/verifier.py}. The verifier confirms weekly coverage (C1), exact group size $s$ (C2), and unique pairs across all weeks (C3) in $O(w g s^2)$ time before setting \texttt{verified = True}.
\end{enumerate}

\section{Computational Experiments}
\label{sec:experiments}

\subsection{Benchmark Suite and Experimental Setup}
\label{sec:exp_setup}

The source list contains 66 triples in the $g$-$p$-$w$ convention and
records the status reported by Kieu et al.~\cite{kieu2025nse}.
We use $s=p$ and pass the triples to the encoder as $(w,g,s)$.
Source statuses serve as cross-checks for decided runs. An unknown
reference status, timeout, or unattempted run provides no decision.

\paragraph{Four-Paradigm Benchmark Setup and Environment.}
For our main comparison (\texttt{benchmarks/results.csv}), we test 64 active instances from \texttt{benchmarks/benchmark\_instances.csv}; the two instances with the largest estimated C3 clause counts, \texttt{11-11-12} ($n=121$) and \texttt{13-13-14} ($n=169$), are turned off in the shared suite for separate CP testing. Each active instance is run across 12 configurations (768 runs total): four SAT C2 encodings $\times$ two SAT solvers (CaDiCaL~1.5.3~\cite{cadical153} and Glucose~4.1~\cite{audemard2018glucose}), Order Encoding $\times$ the same two SAT solvers, Google OR-Tools CP-SAT~\cite{ortools}, and Gurobi ILP~\cite{gurobi}. All 768 runs share the same baseline saved in \texttt{results.csv.meta.json}: without optional \texttt{SB1}/\texttt{SB2} (\texttt{SB=none}), single-threaded execution (\texttt{num\_threads=1}), and a 60-second total time limit (\texttt{timeout\_scope=total}) covering model building, solving, decoding, and verification inside isolated worker processes (after library and Gurobi initialization). To avoid running out of memory on huge CNF files, \texttt{sat} and \texttt{order} skip instances where the C3 clause lower bound $\binom{n}{2}\binom{w}{2}g^2$ exceeds 5{,}000{,}000 clauses (\texttt{SKIPPED\_RESOURCE\_LIMIT}); this skips \texttt{9-9-10} (11{,}809{,}800 estimated C3 clauses) and \texttt{19-10-2} (6{,}481{,}755 estimated C3 clauses). All 768 baseline runs, the 640 extended symmetry-breaking runs (\path{benchmarks/results_sb1.csv} with 384 runs and \path{benchmarks/results_sb_full.csv} with 256 runs under the exact same 60-second single-thread protocol), and the standalone CP-SAT stress runs on \texttt{11-11-12} and \texttt{13-13-14} (\path{benchmarks/results_giant_cp.csv}, \texttt{sb=full}, 1 thread, 60\,s solver limit) were executed on an Acer Swift SF314-512 running Microsoft Windows~11 Home (Build~26300), 12th Gen Intel Core~i5-1240P (12 cores, 16 logical processors), and 16\,GB RAM (15.71\,GiB usable), using Python~3.12.10, PySAT~1.9.dev21~\cite{imms2018pysat}, OR-Tools~9.15.6755, Gurobi~14.0.0 (unrestricted Academic license), and CNSPlots~0.8.0.

\subsection{Experiment 1: Ablation of C2 Exactly-$s$ Encodings}
\label{sec:exp1_c2}

Before screening, we compared the four encodings with an independent
partition enumerator on 29 small instances. The enumeration checks
feasibility without building a CNF. Tests with Glucose agreed with its
answers in all 232 combinations of instance, encoding, and first-week
fixing. We also ran the common CLI with both backends on a feasible
instance and an infeasible instance that passes the partner bound.
All four encodings gave the expected decision in both cases.

We first measured the full CNF generated by each encoder on four cases
with no added symmetry clauses. Table~\ref{tab:tv2_sizes} reports auxiliary
variables and all clauses before solver preprocessing, while
Figure~\ref{fig:tv2_footprint} isolates the C2 clause footprint per
group-week equality across the same four cases. NSC uses the fewest
auxiliary variables in all four cases. Sequential Counter emits the
fewest clauses in these cases because C1 and C3 are identical across
encoders and its C2 contribution is smaller. The two size measures can
therefore favor different encodings.

\begin{table}[pos=!htb]
\centering\small
\caption{Measured formula sizes before solver preprocessing. Each cell reports auxiliary variables / total clauses.}
\label{tab:tv2_sizes}
\begin{tabular}{lrrrr}
\toprule
C2 encoding & \texttt{2-2-2} & \texttt{3-3-3} & \texttt{4-4-4} & \texttt{5-3-5} \\
\midrule
NSC & 20 / 116 & 189 / 1{,}809 & 864 / 15{,}344 & 975 / 30{,}900 \\
SeqCounter & 32 / 104 & 324 / 1{,}728 & 1{,}536 / 15{,}040 & 1{,}800 / 30{,}675 \\
Totalizer & 64 / 160 & 522 / 2{,}268 & 2{,}048 / 17{,}888 & 2{,}950 / 35{,}325 \\
Cardinality Network & 80 / 168 & 1{,}620 / 3{,}519 & 3{,}872 / 17{,}808 & 5{,}400 / 35{,}200 \\
\bottomrule
\end{tabular}
\end{table}

\begin{figure}[pos=!htb]
    \centering
    \includesafeimage[width=0.74\textwidth]{tv2_c2_footprint.pdf}
    \caption{C2 clauses per group-week equality in the four size examples,
    plotted with CNSPlots~0.8.0. Instance labels are $g$-$p$-$w$.}
    \label{fig:tv2_footprint}
\end{figure}

\paragraph{SAT screening protocol.}
We considered each of the 66 source instances with four encodings and
two backends (CaDiCaL~1.5.3~\cite{cadical153} and Glucose~4.1~\cite{audemard2018glucose}) through PySAT~1.9.dev21.
The screen ran each configuration once in sequence, using the fixed
generation order C1, C2, C3 and no added symmetry clauses.
We set a 2-second wall budget on the worker and its solver process.
This includes worker startup, CNF construction and serialization,
solving, and verification. The caller's bookkeeping and timeout
cleanup are outside this budget. To bound resource use, configurations
with a C1+C3 lower estimate above 250{,}000 clauses were skipped.
Instances ruled out by the partner bound use the precheck shortcut
without constructing C1--C3. We record precheck UNSAT separately
from solver UNSAT.

\paragraph{Environment and results.}
This early 2-second SAT-only pilot was run separately on a secondary machine prior to the unified 60-second four-paradigm benchmark in Section~\ref{sec:exp_setup}.
The screen used Windows~10 (build 19045), Python~3.12.14,
and an Intel Core i5-8400H processor (eight logical processors, 15.8~GiB
RAM). The solver wrapper used its backend defaults. Table~\ref{tab:exp1_c2}
reports all 528 configurations. The resource policy skipped
264 configurations across 33 instances. Among the
264 remaining configurations, 102 returned SAT with both model
and schedule verified, 56 returned precheck UNSAT, and 106
timed out. No backend returned UNSAT in this screen.
The SAT decisions cover 15 distinct instances, with up to
45 players. All 158 decided configurations agree
with the known source status for their instance.

\begin{table}[pos=!htb]
\centering\small
\caption{SAT screening on the 66 source instances. Precheck denotes partner-bound UNSAT. The UNSAT column counts only decisions by a backend.}
\label{tab:exp1_c2}
\begin{tabular}{llrrrrr}
\toprule
Backend & C2 encoding & SAT & UNSAT & Precheck & Timeout & Skipped \\
\midrule
CaDiCaL & NSC & 13 & 0 & 7 & 13 & 33 \\
CaDiCaL & SeqCounter & 13 & 0 & 7 & 13 & 33 \\
CaDiCaL & Totalizer & 11 & 0 & 7 & 15 & 33 \\
CaDiCaL & Cardinality Network & 10 & 0 & 7 & 16 & 33 \\
\midrule
Glucose & NSC & 15 & 0 & 7 & 11 & 33 \\
Glucose & SeqCounter & 14 & 0 & 7 & 12 & 33 \\
Glucose & Totalizer & 13 & 0 & 7 & 13 & 33 \\
Glucose & Cardinality Network & 13 & 0 & 7 & 13 & 33 \\
\bottomrule
\end{tabular}
\end{table}

\paragraph{Interpretation.}
This screen measures coverage under the declared budget and resource
policy. The instances reaching each solver and the short cutoff limit
what can be inferred about relative performance. CNF size alone cannot
predict search cost. The source study uses different encodings and
symmetry settings, a different timing metric, and a much longer budget.
We do not use its times as a baseline for speedups.

\paragraph{60-second SAT benchmark protocol.}
The shared benchmark enables 64 of the 66 source instances. The
\texttt{11-11-12} and \texttt{13-13-14} instances are disabled in this
configuration. For SAT we pair each instance with four C2 encodings
and two PySAT backends, CaDiCaL~1.5.3~\cite{cadical153} and
Glucose~4.1~\cite{audemard2018glucose}. We keep C1, C3, and the
generation order fixed and add no symmetry clauses. The configuration
records one thread, a 60-second total wall-time budget, and one result per
configuration. The timer starts after benchmark-worker preparation and
covers CNF construction, the nested SAT subprocess, and verification.
The CSV separately records the backend's solve-only time when available.
Although the metadata file itself omits host details, all runs used the hardware and package environment reported in Section~\ref{sec:exp_setup}.

The partner-bound precheck runs before CNF construction. The benchmark
also skips configurations with more than 5{,}000{,}000 estimated C3
clauses, a lower bound on full CNF size. This affected both backends
and all four encodings on \texttt{9-9-10} and \texttt{19-10-2}.

\paragraph{Coverage and timings.}
Table~\ref{tab:c2-60s-outcomes} summarizes the 512 SAT configurations. Among
them, 240 passed model and schedule verification, 56 were ruled out
by the partner bound, 200 timed out, and 16 were skipped. No backend returned
UNSAT. All 296 decided configurations agree with their reference
statuses. The SAT results cover 33 distinct instances, including
\texttt{17-9-2} with 153 players. On that instance, both backends
solved the Totalizer CNF with 4{,}269{,}414 clauses. C3 alone accounts
for 3{,}360{,}492 clauses in each encoding. With Sequential Counter it
accounts for about 94\% of the total.

\begin{table}[pos=!htb]
\centering\small
\caption{Outcomes on 64 instances per backend and encoding. Precheck denotes partner-bound UNSAT. Median total wall time is computed on the cases solved by all four encodings for the given backend.}
\label{tab:c2-60s-outcomes}
\begin{tabular}{llrrrrr}
\toprule
Backend & C2 encoding & SAT & Precheck & Timeout & Skipped & Median (s) \\
\midrule
CaDiCaL & NSC & 30 & 7 & 25 & 2 & 0.846 \\
CaDiCaL & SeqCounter & 30 & 7 & 25 & 2 & 1.016 \\
CaDiCaL & Totalizer & 30 & 7 & 25 & 2 & 1.412 \\
CaDiCaL & Cardinality Network & 30 & 7 & 25 & 2 & 1.352 \\
\midrule
Glucose & NSC & 32 & 7 & 23 & 2 & 0.590 \\
Glucose & SeqCounter & 31 & 7 & 24 & 2 & 0.574 \\
Glucose & Totalizer & 31 & 7 & 24 & 2 & 0.659 \\
Glucose & Cardinality Network & 26 & 7 & 29 & 2 & 0.713 \\
\bottomrule
\end{tabular}
\end{table}

All four encodings solved the same 30 instances with CaDiCaL and
shared 26 solved instances with Glucose. NSC had the
lowest median total time on the CaDiCaL set and the highest SAT count
with Glucose. Sequential Counter had the lowest median total time on
the Glucose set. Figure~\ref{fig:c2-benchmark-runtime} shows marked
variation by instance size for CaDiCaL. At a fixed number of players,
the group dimensions and the number of weeks can vary. These
measurements describe a single run per configuration, conditional on the
time and resource limits. The timing summaries cover only instances
solved by all four encodings with the selected backend.

\begin{figure}[pos=!htb]
    \centering
    \includesafeimage[width=0.86\textwidth]{runtime_c2_encodings.pdf}
    \caption{Mean total wall time by golfer count for the 30 instances
    solved by all four C2 encodings with CaDiCaL (\texttt{SB=none}) over 17 distinct counts.
    Counts $n\in\{15, 32, 36\}$ aggregate 6, 6, and 4 instances with different $(g,s,w)$ values, whereas sharp dips such as $n=28, 42, 45, 91$ correspond to single low-$w$ instances (\texttt{7-4-2}, \texttt{7-6-3}, \texttt{9-5-2}, \texttt{13-7-2}).}
    \label{fig:c2-benchmark-runtime}
\end{figure}

\paragraph{Runtime breakdown.}
Figure~\ref{fig:c2-benchmark-breakdown} compares the three largest
instances by player count for which CaDiCaL solved all four encodings
and recorded complete phase timings. We subtract CNF construction and
solve-only time from total wall time to obtain the remaining time.
This includes CNF serialization and loading, startup of the SAT
subprocess, and model and schedule verification.

On \texttt{17-9-2}, Sequential Counter emitted 3{,}578{,}670 clauses
and took 27.33 seconds overall. NSC emitted 3{,}583{,}056 clauses and
took 12.82 seconds. The solve phase accounts for most of the gap,
with 15.11 seconds for Sequential Counter and 0.33 seconds for NSC.
For NSC, the remaining time was 9.73 seconds, about 76\% of its total.

\begin{figure}[pos=!htb]
    \centering
    \includesafeimage[width=0.92\textwidth]{runtime_c2_breakdown.pdf}
    \caption{Runtime breakdown for four C2 encodings on the three
    largest commonly solved CaDiCaL instances. Each bar is one run
    with \texttt{SB=none} and a 60-second total budget.}
    \label{fig:c2-benchmark-breakdown}
\end{figure}

\subsection{Experiment 2: Comparing SAT, CP-SAT, and ILP}
\label{sec:exp2_paradigms}

Across all 768 runs in the 60-second benchmark, 330 runs returned \texttt{SAT} with schedules verified by \texttt{src/\allowbreak verifier.py} (covering 34 distinct instances), 84 runs (7 instances) ended early via \texttt{UNSAT\_\allowbreak PRECHECK}, 20 CNF runs (2 instances) hit \texttt{SKIPPED\_\allowbreak RESOURCE\_\allowbreak LIMIT}, and 334 runs reached \texttt{TIMEOUT}, with zero errors (\texttt{ERROR=0}) and zero mismatches against the expected source statuses. Table~\ref{tab:overall_benchmark} summarizes all 12 configurations across the 64 active instances.

\begin{table}[pos=!htb]
\centering\small
\setlength{\tabcolsep}{4pt}
\caption{Summary of all 12 solver configurations across 64 active SGP instances (\texttt{SB=none}, 1 thread, 60\,s total time limit). \textbf{Pre.} is partner-bound \texttt{UNSAT\_PRECHECK}, \textbf{TO} is \texttt{TIMEOUT}, and \textbf{Skip} is the 5M C3 clause guard. Mean, median, and max runtimes are over each row's own solved \texttt{SAT} subset (so rows with different SAT counts are not directly comparable without paired subsets; see Tables~\ref{tab:c2-60s-outcomes} and~\ref{tab:sb_ablation}).}
\label{tab:overall_benchmark}
\begin{tabular}{lllrrrrrrrr}
\toprule
\textbf{Method} & \textbf{Encoding} & \textbf{Solver} & \textbf{SAT} & \textbf{Pre.} & \textbf{TO} & \textbf{Skip} & \textbf{Mean (s)} & \textbf{Med.\ (s)} & \textbf{Max (s)} & \textbf{Max $n$} \\
\midrule
SAT & NSC & CaDiCaL~1.5.3 & 30 & 7 & 25 & 2 & 2.41 & 0.85 & 16.18 & 153 \\
SAT & NSC & Glucose~4.1 & \textbf{32} & 7 & 23 & 2 & 3.04 & 0.86 & 28.46 & 153 \\
SAT & SeqCounter & CaDiCaL~1.5.3 & 30 & 7 & 25 & 2 & 3.98 & 1.02 & 27.33 & 153 \\
SAT & SeqCounter & Glucose~4.1 & 31 & 7 & 24 & 2 & 3.14 & 0.87 & 19.88 & 153 \\
SAT & Totalizer & CaDiCaL~1.5.3 & 30 & 7 & 25 & 2 & 3.21 & 1.41 & 14.82 & 153 \\
SAT & Totalizer & Glucose~4.1 & 31 & 7 & 24 & 2 & 3.30 & 0.99 & 22.44 & 153 \\
SAT & CardNet & CaDiCaL~1.5.3 & 30 & 7 & 25 & 2 & 3.90 & 1.35 & 23.13 & 153 \\
SAT & CardNet & Glucose~4.1 & 26 & 7 & 29 & 2 & 3.13 & 0.71 & 41.13 & 91 \\
\midrule
Order & Order & CaDiCaL~1.5.3 & 20 & 7 & 35 & 2 & 4.72 & 0.87 & 41.04 & 66 \\
Order & Order & Glucose~4.1 & 15 & 7 & 40 & 2 & 4.23 & 0.79 & 29.06 & 66 \\
CP & --- & OR-Tools CP-SAT & 31 & 7 & 26 & 0 & 9.65 & 3.80 & 55.41 & \textbf{190} \\
ILP & --- & Gurobi~14.0 & 24 & 7 & 33 & 0 & 4.00 & 1.09 & 43.67 & 153 \\
\bottomrule
\end{tabular}
\end{table}

Because the unconditioned averages in Table~\ref{tab:overall_benchmark} are computed over different solved sets, raw mean and median values exhibit survivorship bias: for example, \texttt{CardNet}+Glucose reports the lowest unconditioned median (0.71\,s) because its average is restricted to the 26 instances it solved (whereas on those same 26 instances in Table~\ref{tab:c2-60s-outcomes}, it is the slowest of the four encodings at 0.713\,s vs.\ 0.574--0.659\,s), and \texttt{NSC}+Glucose averages 2.02\,s (median 0.75\,s) on the 30 instances solved by \texttt{NSC}+CaDiCaL (vs.\ 2.41\,s, median 0.85\,s), with its higher 3.04\,s unconditioned mean coming solely from solving two additional hard cases (\texttt{8-3-10} in 28.46\,s and \texttt{6-5-5} in 8.09\,s) where CaDiCaL times out.

Figure~\ref{fig:cactus_sat_cp_ilp} compares representative SAT (\texttt{NSC} with CaDiCaL~1.5.3), CP (OR-Tools CP-SAT), and ILP (Gurobi~14.0) runs over the 55 shared attempted instances (excluding the 7 partner-bound prechecks and the 2 CNF-skipped instances). On these 55 instances, SAT (\texttt{NSC}+CaDiCaL) and CP-SAT each solve 30 instances while Gurobi ILP solves 24. Comparing paired runtimes shows clear differences in how the three solvers scale:
\begin{itemize}[noitemsep]
    \item \textbf{Common 24 instances (SAT + CP-SAT + ILP):} By total wall time, SAT (\texttt{NSC}+CaDiCaL) averages 1.68\,s (median 0.53\,s, max 12.82\,s), compared to 4.00\,s for Gurobi ILP (median 1.09\,s, max 43.67\,s) and 4.20\,s for CP-SAT (median 2.43\,s, max 24.16\,s). Gurobi stays below SAT for the first eight points on the cactus curve ($k=1\ldots 8$, up to 0.311\,s vs.\ 0.325\,s at $k=8$) and records a lower total wall time than \texttt{NSC}+CaDiCaL on 12 of the 24 shared instances (including tiny cases like \texttt{3-2-2} at $<$0.01\,s and \texttt{4-3-3} at 0.05\,s, as well as \texttt{8-4-6} and \texttt{9-4-7}). Note that total wall time includes runner overhead that differs by backend: CP-SAT and Gurobi run in-process inside the benchmark worker, whereas \texttt{src/sat\_runner.py} writes the CNF to JSON, spawns a nested Python subprocess, loads the clauses into PySAT before starting \texttt{solver\_wall\_time\_seconds}, and re-checks clauses in Python. In recorded build plus solve time (\texttt{encode\_time\_seconds} $+$ \texttt{solver\_wall\_time\_seconds}, where SAT's solve timer starts after CNF loading), \texttt{NSC}+CaDiCaL is lower than both Gurobi and CP-SAT on all 24/24 shared instances (for example, 0.37\,s vs.\ 1.66\,s on \texttt{8-4-6} and 0.92\,s vs.\ 3.47\,s on \texttt{9-4-7}).
    \item \textbf{Common 30 instances (SAT + CP-SAT):} SAT (\texttt{NSC}+CaDiCaL) averages 2.41\,s (median 0.85\,s) compared to 8.76\,s for CP-SAT (median 3.41\,s), which is a $3.6\times$ speedup on the cases both solvers finish.
    \item \textbf{High-$w$ / near-$w_{\max}$ instances:} As shown in Table~\ref{tab:representative_instances}, the four SAT configurations listed there are faster than CP-SAT and ILP on dense multi-week schedules such as \texttt{5-3-7} (0.98--2.10\,s vs.\ 45.42\,s for CP-SAT and ILP timeout), \texttt{5-5-6} (0.60--1.84\,s vs.\ 55.41\,s for CP-SAT and ILP timeout), and \texttt{9-4-8} (5.56--16.18\,s vs.\ 18.64\,s for CP-SAT, whereas \texttt{CardNet}+Glucose takes 41.13\,s). Moreover, \texttt{8-3-10} (28.46\,s), \texttt{6-5-5} (8.09\,s), and \texttt{9-3-11} (13.08--19.88\,s) are solved only by SAT (Glucose~4.1), while CaDiCaL, Order Encoding, CP-SAT, and Gurobi all time out.
    \item \textbf{Large-$n$, low-$w$ instances (\texttt{19-10-2}, $n=190$):} On \texttt{19-10-2}, the C3 clause lower bound reaches 6{,}481{,}755 clauses, so every SAT and Order run hits \texttt{SKIPPED\_RESOURCE\_LIMIT}, while Gurobi times out after 60\,s. Because \texttt{AllDifferent} and \texttt{Inverse} avoid creating pairwise clauses up front, OR-Tools CP-SAT needs only 37{,}050 variables and 90{,}501 constraints and is the only solver to solve \texttt{19-10-2} (in 36.10\,s). On the other hand, on \texttt{17-9-2} ($n=153$), \texttt{NSC}+CaDiCaL (12.82\,s; 50{,}490 vars, 3{,}583{,}056 clauses) is still nearly $2\times$ faster than CP-SAT (24.16\,s; 24{,}174 vars, 58{,}722 constraints) and $3.4\times$ faster than Gurobi (43.67\,s; 400{,}554 vars, 1{,}198{,}024 constraints).
    \item \textbf{Unproven UNSAT instances:} All eight UNSAT instances that pass the partner-bound precheck (\texttt{5-7-5}, \texttt{6-6-7}, \texttt{5-8-3}, \texttt{6-7-4}, \texttt{5-10-4}, \texttt{5-10-5}, \texttt{6-9-4}, and \texttt{8-10-3}) timed out at 60\,s across all 12 configurations under \texttt{SB=none}. As shown in Section~\ref{sec:exp3_order_sb}, with \texttt{sb1}, all six tested configurations proved seven of the eight remaining UNSAT instances within the 60-second budget (and OR-Tools CP-SAT proved all eight under \texttt{full}).
\end{itemize}

\begin{figure}[pos=!htb]
    \centering
    \includesafeimage[width=0.88\textwidth]{cactus_sat_cp_ilp.pdf}
    \caption{Cactus plot comparing SAT (\texttt{NSC} + CaDiCaL~1.5.3, 30 solved), CP (OR-Tools CP-SAT, 30 solved), and ILP (Gurobi~14.0, 24 solved) over 55 shared attempted instances (\texttt{SB=none}, 1 thread, 60\,s total limit), plotted with CNSPlots~0.8.0. Each curve sorts the solved runtimes of its own solver independently.}
    \label{fig:cactus_sat_cp_ilp}
\end{figure}

\begin{table}[pos=!htb]
\centering\footnotesize
\setlength{\tabcolsep}{3.5pt}
\caption{Total wall-time runtime (seconds) on 19 representative SGP instances ($n \in [6, 190]$, \texttt{SB=none}, 60\,s limit). \textbf{TO} = timeout; \textbf{SKIP} = exceeded 5M C3 clause guard. Fastest total wall time among the configurations shown is in bold (note that SAT/Order totals include nested subprocess/JSON IPC overhead).}
\label{tab:representative_instances}
\begin{tabular}{lrccccccc}
\toprule
\textbf{Inst.} & $n$ & \textbf{NSC(CaD)} & \textbf{NSC(Glu)} & \textbf{Seq(Glu)} & \textbf{Tot(Glu)} & \textbf{Order(CaD)} & \textbf{CP-SAT} & \textbf{Gurobi} \\
\midrule
\texttt{3-2-2} & 6 & 0.19 & 0.20 & 0.21 & 0.22 & 0.22 & 0.03 & \textbf{$<$0.01} \\
\texttt{5-3-6} & 15 & 0.52 & 0.39 & \textbf{0.38} & 0.45 & 8.41 & 1.75 & 2.08 \\
\texttt{5-3-7} & 15 & \textbf{0.98} & 1.35 & 2.10 & 1.01 & TO & 45.42 & TO \\
\texttt{5-4-4} & 20 & 0.32 & 0.32 & \textbf{0.31} & 0.35 & 1.01 & 0.79 & 1.65 \\
\texttt{8-3-10} & 24 & TO & \textbf{28.46} & TO & TO & TO & TO & TO \\
\texttt{5-5-6} & 25 & 1.84 & 0.62 & \textbf{0.60} & 0.70 & 0.94 & 55.41 & TO \\
\texttt{9-3-11} & 27 & TO & TO & 19.88 & \textbf{13.08} & TO & TO & TO \\
\texttt{6-5-5} & 30 & TO & \textbf{8.09} & TO & TO & TO & TO & TO \\
\texttt{8-4-6} & 32 & 1.97 & 1.92 & 1.94 & 2.03 & TO & 4.10 & \textbf{1.68} \\
\texttt{8-4-7} & 32 & 4.51 & \textbf{3.39} & 4.65 & 5.53 & TO & 14.74 & TO \\
\texttt{9-4-7} & 36 & 4.11 & 4.00 & 4.03 & 4.40 & TO & 6.23 & \textbf{3.51} \\
\texttt{9-4-8} & 36 & 16.18 & \textbf{5.56} & 10.40 & 14.76 & TO & 18.64 & TO \\
\texttt{7-7-3} & 49 & \textbf{0.83} & \textbf{0.83} & 0.87 & 0.99 & 1.70 & 5.25 & 1.93 \\
\texttt{11-6-2} & 66 & 1.30 & \textbf{1.23} & 1.45 & 1.48 & 2.88 & 6.06 & 2.60 \\
\texttt{9-9-3} & 81 & 3.67 & \textbf{3.06} & 4.36 & 4.90 & TO & 11.34 & TO \\
\texttt{10-9-3} & 90 & 4.95 & \textbf{4.75} & 7.72 & 6.39 & TO & 16.67 & TO \\
\texttt{15-8-2} & 120 & \textbf{6.11} & 6.14 & 6.91 & 9.73 & TO & 17.81 & 21.45 \\
\texttt{17-9-2} & 153 & \textbf{12.82} & 13.07 & 18.49 & 22.44 & TO & 24.16 & 43.67 \\
\texttt{19-10-2} & 190 & SKIP & SKIP & SKIP & SKIP & SKIP & \textbf{36.10} & TO \\
\bottomrule
\end{tabular}
\end{table}

\paragraph{Standalone CP-SAT Runs on High-Clause Disabled Instances.}
To test model construction and search beyond the 64-instance suite, we ran OR-Tools CP-SAT alone (\texttt{sb=full}, 1 thread, 60\,s solver limit) on the two disabled instances with the largest estimated C3 clause counts, \texttt{11-11-12} ($n=121$) and \texttt{13-13-14} ($n=169$). Both inputs lie at the partner bound $w = w_{\max}$ (12 and 14 weeks), and their estimated C3 clause counts reach 57.98\,M and 218.32\,M clauses, exceeding the benchmark's 5M-clause skip threshold for SAT and Order Encoding. As summarized in Table~\ref{tab:giant_cp}, CP-SAT builds the two models in 1.29\,s (91{,}476 variables, 184{,}538 constraints) and 2.85\,s (205{,}842 variables, 416{,}600 constraints), and both runs then reach \texttt{TIMEOUT} at the 60\,s solver limit (for reference, Kieu et al.~\cite{kieu2025nse} report 146\,s on \texttt{11-11-12} and a timeout on \texttt{13-13-14}).

\begin{table}[pos=!htb]
\centering\small
\setlength{\tabcolsep}{4pt}
\caption{Standalone OR-Tools CP-SAT runs on the two disabled instances with the largest estimated C3 clause counts (\texttt{sb=full}, 1 thread, 60\,s solver limit). Est.\ C3 gives the theoretical CNF C3 clause count that exceeds the 5M-clause benchmark skip rule for SAT/Order.}
\label{tab:giant_cp}
\begin{tabular}{lrrrrrrr}
\toprule
\textbf{Inst.} & $n$ ($w = w_{\max}$) & \textbf{Est.\ C3 Clauses} & \textbf{CP Vars} & \textbf{CP Constr.} & \textbf{Build (s)} & \textbf{Solve (s)} & \textbf{Total (s) / Status} \\
\midrule
\texttt{11-11-12} & 121 (12) & 57{,}978{,}360 & 91{,}476 & 184{,}538 & 1.29 & 60.12 & 61.48 (\texttt{TIMEOUT}) \\
\texttt{13-13-14} & 169 (14) & 218{,}320{,}284 & 205{,}842 & 416{,}600 & 2.85 & 61.75 & 64.66 (\texttt{TIMEOUT}) \\
\bottomrule
\end{tabular}
\end{table}

\subsection{Experiment 3: Direct vs.\ Order Encoding and Symmetry-Breaking Ablation}
\label{sec:exp3_order_sb}

Figure~\ref{fig:bar_direct_vs_order} compares three-index Direct SAT (\texttt{NSC}) with slot-based Order Encoding under the same CaDiCaL~1.5.3, single-thread, and \texttt{SB=none} settings across all 20 commonly solved instances ordered by $n$. Across all 64 instances, Direct (\texttt{NSC}) solves 30 instances with CaDiCaL (32 with Glucose) up to $n=153$, whereas Order Encoding solves 20 with CaDiCaL (15 with Glucose) up to $n=66$ (\texttt{11-6-2}) and solves no instance with $n > 66$ (timing out on all attempted instances with $n \ge 80$, including \texttt{8-10-3} at $n=80$, while \texttt{9-9-10} and \texttt{19-10-2} are skipped by the 5M C3 clause guard).

\begin{figure}[pos=!htb]
    \centering
    \includesafeimage[width=0.88\textwidth]{bar_direct_vs_order.pdf}
    \caption{Total runtime comparison between Direct (\texttt{NSC}) and Order Encoding on CaDiCaL~1.5.3 (\texttt{SB=none}, 60\,s limit) across all 20 shared solved instances, shown as a log-scale dot plot with CNSPlots~0.8.0. Across these 20 instances, the $6.6\times$ speedup refers to the ratio of mean total runtimes (4.72\,s vs.\ 0.71\,s).}
    \label{fig:bar_direct_vs_order}
\end{figure}

Across all 20 instances solved by both encodings with CaDiCaL, Direct (\texttt{NSC}) takes 0.71\,s on average using 3{,}574 variables and 131{,}909 clauses, while Order Encoding takes 4.72\,s using 44{,}426 variables and 207{,}932 clauses---so Direct uses $12.4\times$ fewer variables and achieves a $6.6\times$ lower mean runtime across the 20 shared cases (4.72\,s vs.\ 0.71\,s). Order Encoding slows down sharply as $n$ and $w$ grow, such as on \texttt{5-3-6} (0.52\,s vs.\ 8.41\,s), \texttt{7-6-3} (0.68\,s vs.\ 27.83\,s), and \texttt{9-4-6} (2.91\,s vs.\ 41.04\,s), because it creates $O(w g s n)$ order and channeling variables plus $\Theta(w n^3)$ sequential-counter \texttt{AllDifferent} auxiliary variables and clauses. Among the 20 shared CaDiCaL cases, \texttt{5-5-6} is the only instance where Order Encoding achieves a genuine solving speedup (0.94\,s total / 0.11\,s solve vs.\ 1.84\,s total / 1.21\,s solve for \texttt{NSC}+CaDiCaL, though \texttt{NSC}+Glucose finishes in 0.62\,s); their total times were nearly equal on \texttt{5-3-4} (0.285\,s for Order vs.\ 0.290\,s for Direct), while Direct recorded a lower build plus solve time (0.010\,s vs.\ 0.018\,s).

\paragraph{Impact of Static Symmetry Breaking (\texttt{none} vs.\ \texttt{sb1} vs.\ \texttt{full}).}
To measure the effect of the static symmetry-breaking modes from Section~\ref{sec:symmetry_breaking}, we ran 640 additional runs across all 64 active instances under the exact same single-thread, 60-second total budget: 384 \texttt{sb1} runs covering six representative configurations (\path{benchmarks/results_sb1.csv}) and 256 \texttt{full} runs covering the four configurations that support \texttt{SB2} (\path{benchmarks/results_sb_full.csv}). Every run finished with zero errors (\texttt{ERROR=0}) and zero mismatches against reference statuses; together with \texttt{SB=none}, the solvers decide 50 of the 64 active instances overall (35 distinct SAT, 8 solver-proved UNSAT, and 7 partner-bound precheck). Table~\ref{tab:sb_ablation} and Figure~\ref{fig:bar_sb} summarize the results across all three modes:

\begin{table}[pos=!htb]
\centering\small
\setlength{\tabcolsep}{4pt}
\caption{Comparison across \texttt{SB=none}, \texttt{sb1}, and \texttt{full} (\texttt{SB1}+\texttt{SB2}) on 64 active SGP instances (1 thread, 60\,s total limit). Cells show \textbf{SAT / UNSAT / TO} excluding 7 \texttt{UNSAT\_PRECHECK} runs (and 2 C3 clause-guard skips on SAT/Order). Order Encoding does not implement \texttt{SB2}. Paired mean time averages total runtime across commonly solved SAT runs ($n_{\mathrm{sh}}$).}
\label{tab:sb_ablation}
\begin{tabular}{lccccc}
\toprule
\textbf{Method / Solver} & \textbf{none} & \textbf{sb1} & \textbf{full} & \textbf{Paired Mean SAT Time (s)} & $n_{\mathrm{sh}}$ \\
\midrule
SAT NSC (CaDiCaL) & 30 / 0 / 25 & 30 / 7 / 18 & 31 / 7 / 17 & 2.41 $\to$ 2.39 $\to$ \textbf{2.24} & 30 \\
SAT NSC (Glucose) & 32 / 0 / 23 & 31 / 7 / 17 & 32 / 7 / \textbf{16} & 2.02 $\to$ 2.05 $\to$ \textbf{2.01} & 30 \\
Order (CaDiCaL) & 20 / 0 / 35 & 25 / 7 / 23 & --- & 4.72 $\to$ \textbf{1.54} & 20 \\
Order (Glucose) & 15 / 0 / 40 & 20 / 7 / 28 & --- & 4.23 $\to$ \textbf{0.89} & 15 \\
OR-Tools CP-SAT & 31 / 0 / 26 & 30 / 7 / 20 & \textbf{32} / \textbf{8} / 17 & 9.35 $\to$ 4.29 $\to$ \textbf{3.34} & 30 \\
Gurobi 14.0 ILP & 24 / 0 / 33 & 28 / 7 / 22 & 28 / 7 / 22 & 4.00 $\to$ 4.80 $\to$ \textbf{3.24} & 24 \\
\bottomrule
\end{tabular}
\end{table}

\begin{figure}[pos=!htb]
    \centering
    \includesafeimage[width=0.98\textwidth]{bar_symmetry_breaking.pdf}
    \caption{Effect of static symmetry breaking (\texttt{none}, \texttt{sb1}, \texttt{full}) across 64 active instances (1 thread, 60\,s total limit): (a) solver-decided SAT + UNSAT counts excluding precheck; (b) paired mean total runtime on shared solved SAT instances (with each solver's shared sample size $n_{\mathrm{sh}}$ shown beneath the axis labels), plotted with CNSPlots~0.8.0.}
    \label{fig:bar_sb}
\end{figure}

\begin{itemize}[noitemsep]
    \item \textbf{UNSAT breakthrough ($0/8 \to 7/8 \to 8/8$):} Under \texttt{SB=none}, none of the configurations proved any of the 8 non-precheck UNSAT instances within 60\,s. Once \texttt{sb1} fixes week~1, all six configurations immediately prove 7 of them (\texttt{5-7-5}, \texttt{5-8-3}, \texttt{6-7-4}, \texttt{5-10-4}, \texttt{5-10-5}, \texttt{6-9-4}, \texttt{8-10-3}) in 0.13--5.44\,s (all seven satisfy $s > g$, triggering immediate week-2 pigeonhole conflicts). Under \texttt{full}, OR-Tools CP-SAT is the only solver to also prove the tight symmetric case \texttt{6-6-7} ($g=s=6$) in 7.05\,s, deciding all 15 UNSAT reference instances (8 by solver + 7 by precheck).
    \item \textbf{CP-SAT gains and speedup:} With \texttt{full}, CP-SAT drops its paired mean SAT runtime by $2.8\times$ (9.35\,s $\to$ 3.34\,s), cuts \texttt{19-10-2} from 36.10\,s to 29.80\,s (and 20.26\,s under \texttt{sb1}), and unlocks \texttt{7-7-8} in 15.08\,s (never solved under \texttt{SB=none}), reaching 47 decided instances (32 SAT + 8 UNSAT + 7 precheck); under \texttt{sb1} alone, CP-SAT temporarily loses \texttt{9-4-8} (31 $\to$ 30 SAT) before recovering it under \texttt{full} (12.70\,s).
    \item \textbf{Order Encoding and Gurobi ILP gains:} Fixing week~1 helps Order Encoding the most on runtime and SAT coverage: both CaDiCaL and Glucose add $+5$ SAT (reaching \texttt{13-7-2} at $n=91$) and cut mean runtime by $3.1\times$ (4.72\,s $\to$ 1.54\,s) and $4.8\times$ (4.23\,s $\to$ 0.89\,s). Gurobi ILP adds $+4$ SAT instances under \texttt{sb1}/\texttt{full} (\texttt{5-5-6}, \texttt{9-9-3}, \texttt{10-9-3}, and even \texttt{19-10-2} at $n=190$ in 49.47\,s under \texttt{sb1} and 43.70\,s under \texttt{full}).
    \item \textbf{SAT \texttt{NSC} trade-offs:} \texttt{NSC}+Glucose already unlocks \texttt{9-3-11} under \texttt{sb1} (15.22\,s) and speeds up to 8.61\,s under \texttt{full}, whereas \texttt{NSC}+CaDiCaL first unlocks it under \texttt{full} (30.81\,s), leaving Glucose with only 16 timeouts. However, static symmetry constraints also introduce occasional SAT regressions: under \texttt{sb1}, \texttt{NSC}+Glucose times out on both \texttt{8-3-10} and \texttt{6-5-5} (dropping from 32 to 31 SAT) before recovering \texttt{8-3-10} under \texttt{full} (31.06\,s) while still timing out on \texttt{6-5-5} (vs.\ 8.09\,s at \texttt{none}), showing that static symmetry clauses can occasionally alter CDCL branching paths unfavorably on satisfiable inputs.
\end{itemize}

\section{Conclusion}
\label{sec:conclusion}

In this project, we built and tested a shared benchmark pipeline for four exact approaches to the Social Golfer Problem: three-index one-hot SAT with four \textit{Exactly}-$s$ encodings, Order Encoding, OR-Tools CP-SAT, and Gurobi~14.0 ILP. Across 768 single-threaded baseline runs (\texttt{SB=none}) and 640 symmetry-breaking runs (\texttt{sb1} and \texttt{full}) on 64 active instances with a 60-second limit (plus two standalone CP-SAT runs on the disabled high-clause instances), our results highlight four takeaways:
\begin{enumerate}[label=(\arabic*),noitemsep]
    \item Among the C2 encodings, \texttt{NSC} produces the smallest variable count ($5.9\times$ fewer variables than \texttt{Cardinality Network} on the 30 common CaDiCaL runs), gives the lowest average runtime among the four C2 encodings on those 30 CaDiCaL instances (2.41\,s at \texttt{SB=none}), and solves the most instances overall (32 on Glucose~4.1, averaging 2.02\,s on the same 30 common cases).
    \item Three-index Direct SAT outperforms slot-based Order Encoding in both size ($12.4\times$ fewer variables on shared cases) and speed ($6.6\times$ ratio of mean runtimes across the 20 shared CaDiCaL cases at \texttt{SB=none}, reaching $n=153$ instead of $n=66$).
    \item Across the three solver families without optional \texttt{SB1}/\texttt{SB2} (\texttt{SB=none}), SAT handles tightly constrained multi-week instances best (uniquely solving \texttt{8-3-10}, \texttt{6-5-5}, and \texttt{9-3-11}), Gurobi ILP records lower total wall time on small inputs (where SAT's total wall time includes nested subprocess, JSON serialization, and PySAT loading overhead, while \texttt{NSC}+CaDiCaL has lower recorded build plus solve time on all 24 shared instances) before $O(n^2 g w)$ linearization variables slow down search, and OR-Tools CP-SAT scales to large $n$ (\texttt{19-10-2} at $n=190$, and builds \texttt{11-11-12} and \texttt{13-13-14} in under 3\,s) when CNF clauses hit the 5M C3 clause guard.
    \item Within the 60-second budget, static symmetry breaking markedly improves UNSAT coverage and domain-based search: with \texttt{sb1}, all six tested configurations proved seven of the eight remaining UNSAT instances within the 60-second budget (up from 0/8 under \texttt{SB=none}) and Order Encoding sped up by $3.1\times$--$4.8\times$, while \texttt{full} (\texttt{SB1}+\texttt{SB2}) enabled CP-SAT to prove all 8/8 non-precheck UNSAT cases, speed up by $2.8\times$, and unlock \texttt{7-7-8} (deciding 47/64 instances on its own, and bringing the union across all tested configurations to 50/64).
\end{enumerate}

\paragraph{Limitations and Future Work.}
First, with a single-threaded 60\,s cutoff, 14 hard active instances (nine of which are also marked \texttt{UNKNOWN} in Kieu et al.~\cite{kieu2025nse}, plus \texttt{9-9-10}, which hits the 5M C3 clause guard on SAT/Order and times out on CP-SAT/Gurobi) and the two disabled high-clause boundary cases \texttt{11-11-12} and \texttt{13-13-14} remain undecided; multi-threaded search, longer time budgets, static lexicographic group ordering, and dynamic symmetry breaking (such as SBDS/SBDD) are natural next steps. Second, replacing pairwise four-literal C3 clauses with auxiliary pair variables could help CNF encodings avoid hitting the 5M C3 clause guard at large $n$.

\section*{Group Contributions and Acknowledgments}
This project is a collaboration among three students for course INT3011E (Advanced Topics in Computer Science) at VNU University of Engineering and Technology under the supervision of Dr.~To Van Khanh:
\begin{itemize}[noitemsep]
    \item \textbf{Do Huy Hieu (Member 1)}: Universal Solution Verifier, schedule display, shared CLI and DIMACS exporter, OR-Tools CP-SAT model, \texttt{cnsplots} visualization, full benchmark execution (768 baseline runs, 640 symmetry-breaking runs, and giant CP-SAT tests), and paper editing (Abstract, Introduction, CP formulation, Pipeline, Experiments, Conclusion).
    \item \textbf{Pham Van Huong (Member 2)}: Three-index one-hot SAT encoding, custom NSC and PySAT \textit{Exactly}-$s$ encodings, isolated PySAT solver runner, CSV~v1 benchmark schema, dependencies, SAT methodology, the SAT-only screening experiments, and the C2 benchmark analysis.
    \item \textbf{Nguyen Minh Son (Member 3)}: Order Encoding, ILP formulation, symmetry-breaking module, automated benchmark suite, and ILP/Order/Symmetry sections.
\end{itemize}
We thank Dr.~To Van Khanh for his guidance, lectures, and reference papers on SGP and Order Encoding.

\section*{Data and Code Availability}
The solver source code, instance lists, 60-second benchmark CSV logs, and the plotting script for Figures~\ref{fig:c2-benchmark-runtime}--\ref{fig:bar_sb} are publicly available in the project GitHub repository at \url{https://github.com/huong1105/INT3011E-SGP} (with the initial \texttt{SB=none} baseline snapshot archived at commit \texttt{aa79113}, and the extended \texttt{sb1}/\texttt{full}, high-clause CP-SAT, and updated figure files hosted on the \texttt{main} branch).
Specifically, the reference table and instance lists are \path{data/sgp_paper_instances.csv}, \path{benchmarks/benchmark_instances.csv}, \path{data/giant_11-11-12.csv}, and \path{data/giant_13-13-14.csv}; the result CSV files (with their \path{.meta.json} manifests) are \path{benchmarks/results.csv}, \path{benchmarks/results_sb1.csv}, \path{benchmarks/results_sb_full.csv}, and \path{benchmarks/results_giant_cp.csv}.
The SAT pilot and automated four-paradigm runners are \path{tools/run_sat_benchmark.py} and \path{benchmarks/run_all.py}, and the figure-generation script is \path{benchmarks/plot_cns.py}, which builds Figures~\ref{fig:c2-benchmark-runtime}, \ref{fig:c2-benchmark-breakdown}, \ref{fig:cactus_sat_cp_ilp}, \ref{fig:bar_direct_vs_order}, and~\ref{fig:bar_sb} from the saved CSV files.

\begin{thebibliography}{99}
\small
\setlength{\itemsep}{1pt plus 0.5pt}
\setlength{\parskip}{0pt}

\bibitem{asin2009cardnet}
R.~As{\'\i}n, R.~Nieuwenhuis, A.~Oliveras, and E.~Rodr{\'\i}guez-Carbonell,
``Cardinality Networks and Their Applications,''
in \textit{Theory and Applications of Satisfiability Testing (SAT 2009)},
LNCS, vol.~5584, pp.~167--180, Springer, 2009.

\bibitem{audemard2018glucose}
G.~Audemard and L.~Simon,
``On the Glucose SAT Solver,''
\textit{International Journal on Artificial Intelligence Tools}, vol.~27, no.~1, 1840001, 2018.

\bibitem{azevedo2006dynamic}
F.~Azevedo,
``An Attempt to Dynamically Break Symmetries in the Social Golfers Problem,''
in \textit{Recent Advances in Constraints (CSCLP 2006)},
LNCS, vol.~4651, pp.~33--47, Springer, 2007.

\bibitem{bailleux2003totalizer}
O.~Bailleux and Y.~Boufkhad,
``Efficient CNF Encoding of Boolean Cardinality Constraints,''
in \textit{Principles and Practice of Constraint Programming (CP 2003)},
LNCS, vol.~2833, pp.~108--122, Springer, 2003.

\bibitem{biere2021handbook}
A.~Biere, M.~Heule, H.~van~Maaren, and T.~Walsh (Eds.),
\textit{Handbook of Satisfiability},
vol.~336 of \textit{Frontiers in Artificial Intelligence and Applications}, 2nd ed., IOS Press, 2021.

\bibitem{cadical153}
A.~Biere et al., \textit{CaDiCaL 1.5.3}, software release.
\url{https://github.com/arminbiere/cadical/releases/tag/rel-1.5.3}

\bibitem{briskorn2008sports}
D.~Briskorn,
\textit{Sports Leagues Scheduling: Models, Combinatorial Properties, and Optimization Algorithms},
Lecture Notes in Economics and Mathematical Systems, vol.~603, Springer, Berlin, 2008.

\bibitem{colbourn2006handbook}
C.~J.~Colbourn and J.~H.~Dinitz (Eds.),
\textit{Handbook of Combinatorial Designs},
2nd ed., CRC Press, Boca Raton, FL, 2006.

\bibitem{cotta2006memetic}
C.~Cotta, I.~Dot{\'u}, A.~J.~Fern{\'a}ndez, and P.~Van~Hentenryck,
``Scheduling Social Golfers with Memetic Evolutionary Programming,''
in \textit{Hybrid Metaheuristics}, LNCS, vol.~4030, pp.~150--161, Springer, 2006.

\bibitem{dotu2007local}
I.~Dot{\'u} and P.~Van~Hentenryck,
``Scheduling Social Tournaments Locally,''
\textit{AI Communications}, vol.~20, no.~3, pp.~151--162, 2007.

\bibitem{gent2005sat}
I.~Gent and I.~Lynce,
``A SAT Encoding for the Social Golfer Problem,''
in \textit{IJCAI Workshop on Modelling and Solving Problems with Constraints}, 2005.

\bibitem{ortools}
Google Developers,
\textit{OR-Tools CP-SAT Solver Reference Documentation},
\url{https://developers.google.com/optimization}, 2024.

\bibitem{gurobi}
Gurobi Optimization, LLC,
\textit{Gurobi Optimizer Reference Manual},
\url{https://www.gurobi.com}, 2025.

\bibitem{harvey2001csplib}
W.~Harvey,
``CSPLib Problem 010: Social Golfers Problem,''
\url{https://www.csplib.org/Problems/prob010/}, 2001.

\bibitem{harvey2005molr}
W.~Harvey and T.~Winterer,
``Solving the MOLR and Social Golfers Problems,''
in \textit{Principles and Practice of Constraint Programming (CP 2005)},
LNCS, vol.~3709, pp.~286--300, Springer, 2005.

\bibitem{henz2001basketball}
M.~Henz,
``Scheduling a Major College Basketball Conference---Revisited,''
\textit{Operations Research}, vol.~49, no.~1, pp.~163--168, 2001.

\bibitem{imms2018pysat}
A.~Ignatiev, A.~Morgado, and J.~Marques-Silva,
``PySAT: A Python Toolkit for Prototyping with SAT Oracles,''
in \textit{Theory and Applications of Satisfiability Testing (SAT 2018)},
LNCS, vol.~10929, pp.~428--437, Springer, 2018.

\bibitem{kieu2025nse}
T.~V.~Kieu, N.~T.~Nguyen, and K.~V.~To,
``An Efficient SAT Encoding for Solving the Social Golfer Problem,''
\textit{RAIRO -- Operations Research}, vol.~60, no.~1, pp.~173--199, 2026. \url{https://doi.org/10.1051/ro/2025155}

\bibitem{lardeux2015sce}
F.~Lardeux, E.~Monfroy, B.~Crawford, and R.~Soto,
``Set Constraint Model and Automated Encoding into SAT: Application to the Social Golfer Problem,''
\textit{Annals of Operations Research}, vol.~235, no.~1, pp.~423--452, 2015.

\bibitem{lester2021mahjong}
M.~M.~Lester,
``Scheduling Reach Mahjong Tournaments Using Pseudoboolean Constraints,''
in \textit{Theory and Applications of Satisfiability Testing (SAT 2021)},
LNCS, vol.~12831, pp.~349--358, Springer, 2021.
\url{https://doi.org/10.1007/978-3-030-80223-3_24}

\bibitem{liu2019sgp}
K.~Liu, S.~L{\"o}ffler, and P.~Hofstedt,
``Social Golfer Problem Revisited,''
in \textit{Agents and Artificial Intelligence (ICAART 2019, Revised Selected Papers)},
LNAI, vol.~11978, pp.~72--99, Springer, 2019.

\bibitem{margot2002pruning}
F.~Margot,
``Pruning by Isomorphism in Branch-and-Cut,''
\textit{Mathematical Programming}, vol.~94, no.~1, pp.~71--90, 2002.

\bibitem{montgomery2017design}
D.~C.~Montgomery,
\textit{Design and Analysis of Experiments},
9th ed., John Wiley \& Sons, Hoboken, NJ, 2017.

\bibitem{pysat-card-docs}
PySAT developers, \textit{Cardinality Encodings (\texttt{pysat.card})}.
\url{https://pysathq.github.io/docs/html/api/card.html}

\bibitem{pysat-network-source}
PySAT developers, \textit{Cardinality Network implementation}.
\url{https://github.com/pysathq/pysat/blob/master/cardenc/sortcard.hh}

\bibitem{pysat-seqcounter-source}
PySAT developers, \textit{Sequential Counter implementation}.
\url{https://github.com/pysathq/pysat/blob/master/cardenc/seqcounter.hh}

\bibitem{pysat-totalizer-source}
PySAT developers, \textit{Totalizer implementation}.
\url{https://github.com/pysathq/pysat/blob/master/cardenc/mto.hh}

\bibitem{sinz2005towards}
C.~Sinz,
``Towards an Optimal CNF Encoding of Boolean Cardinality Constraints,''
in \textit{Principles and Practice of Constraint Programming (CP 2005)},
LNCS, vol.~3709, pp.~827--831, Springer, 2005.

\bibitem{stinson2004designs}
D.~R.~Stinson,
\textit{Combinatorial Designs: Constructions and Analysis},
Springer, New York, 2004.

\bibitem{tamura2009order}
N.~Tamura, A.~Taga, S.~Kitagawa, and M.~Banbara,
``Compiling Finite Linear CSP into SAT,''
\textit{Constraints}, vol.~14, no.~2, pp.~254--272, 2009.

\bibitem{triska2012tme}
M.~Triska and N.~Musliu,
``An Improved SAT Formulation for the Social Golfer Problem,''
\textit{Annals of Operations Research}, vol.~194, no.~1, pp.~427--438, 2012.

\bibitem{verplancke2026}
H.~Verplancke, D.~Van Bulck, V.~Lim\`ere, and N.~A.~Schmid,
``The (a)social exam invigilator assignment problem,''
\textit{Journal of Scheduling}, vol.~29, no.~3, pp.~185--202, 2026.
\url{https://doi.org/10.1007/s10951-025-00868-7}

\bibitem{walsh2006symmetry}
T.~Walsh,
``General Symmetry Breaking Constraints,''
in \textit{Principles and Practice of Constraint Programming (CP 2006)},
LNCS, vol.~4204, pp.~650--664, Springer, 2006.

\end{thebibliography}

\end{document}
